\documentclass[10pt,a4paper]{report}
\usepackage[utf8]{inputenc}
\usepackage[margin=1in]{geometry}
\usepackage{amsmath,amssymb,amsfonts}
\usepackage{booktabs,tabularx}
\usepackage{xcolor}
\usepackage{listings}
\usepackage{makeidx}
\usepackage[colorlinks=true,linkcolor=blue,citecolor=blue,urlcolor=blue]{hyperref}

\makeindex

% Define code syntax highlighting colors
\definecolor{codegreen}{rgb}{0,0.6,0}
\definecolor{codegray}{rgb}{0.5,0.5,0.5}
\definecolor{codepurple}{rgb}{0.58,0,0.82}
\definecolor{backcolour}{rgb}{0.96,0.96,0.96}

\lstdefinestyle{pythonstyle}{
    backgroundcolor=\color{backcolour},   
    commentstyle=\color{codegreen},
    keywordstyle=\color{blue},
    numberstyle=\tiny\color{codegray},
    stringstyle=\color{codepurple},
    basicstyle=\ttfamily\footnotesize,
    breakatwhitespace=false,         
    breaklines=true,                 
    captionpos=b,                    
    keepspaces=true,                 
    numbers=left,                    
    numbersep=5pt,                  
    showspaces=false,                
    showstringspaces=false,
    showtabs=false,                  
    tabsize=4
}
\lstset{style=pythonstyle}

\title{\textbf{VibeVolts Simulation Toolkit}\\
\large Literate Programming Source Documentation}
\author{\textbf{VibeVolts Development Team}}
\date{\today}

\begin{document}

\maketitle

\begin{abstract}
This document presents the complete literate source code of the \textbf{VibeVolts} space environment simulation toolkit. VibeVolts models space-based and ground-based optical sensors, celestial motion, satellite dynamics, target radiometry, and discrete-event scan scheduling. In accordance with literate programming principles, every module's source code is surrounded by mathematical derivations, system architecture explanations, and hyperlinked cross-references.
\end{abstract}

\tableofcontents
\newpage

\chapter{Architecture \& Foundations}
\section{System Overview}
The VibeVolts architecture is built around a centralized state container (\texttt{sim\_data}, an instance of \texttt{SimulationState}). The state encapsulates:
\begin{enumerate}
    \item \textbf{Kinematics}: Positions ($\mathbf{r} \in \mathbb{R}^3$), velocities ($\mathbf{v} \in \mathbb{R}^3$), and accelerations ($\mathbf{a} \in \mathbb{R}^3$) in the Geocentric Celestial Reference System (GCRS).
    \item \textbf{Assets}: Satellites (\texttt{SatellitesState}), Ground Stations (\texttt{ObservatoriesState}), Celestial Bodies (\texttt{CelestialState}), and Stationary Targets (\texttt{FixedPointsState}).
    \item \textbf{Sensors}: Parallelized optoelectronic parameters and pointing vectors (\texttt{DetectorArray}).
    \item \textbf{Discrete Event Scheduler}: Cadence scheduling (\texttt{CadenceState}) organizing detectors by integration time.
\end{enumerate}



\chapter{minimalsimulation.py: Core Simulation State Infrastructure}
\label{sec:minimalsimulation}
\index{minimalsimulation.py}

\section*{Overview}
Defines SimulationState, ComponentState, DotDict, SchemaDict, and get\_detector\_positions().

\section*{Literate Source Code}
The source code of \texttt{minimalsimulation.py} is documented below. Section label: \ref{sec:minimalsimulation}.

\begin{lstlisting}[language=Python,caption={minimalsimulation.py Source Code},label={code:minimalsimulation}]
import numpy as np
from datetime import datetime, timezone
from typing import Dict, Any, Optional, List
from dataclasses import dataclass, field
from constants import *


def safe_normalize(v: np.ndarray, axis: int = -1) -> np.ndarray:
    """
    Safely normalizes an array of vectors along the specified axis,
    preventing division by zero.
    """
    norms = np.linalg.norm(v, axis=axis, keepdims=True)
    return np.where(norms == 0, 0.0, v / norms)


class DotDict(dict):
    """
    A clean dictionary subclass that supports dot-notation attribute access.
    Unifies attribute access and dictionary keys into a single storage mechanism.
    """
    def __getattr__(self, name):
        try:
            return self[name]
        except KeyError:
            raise AttributeError(f"'{self.__class__.__name__}' object has no attribute '{name}'")

    def __setattr__(self, name, value):
        self[name] = value

    def get(self, key, default=None):
        try:
            return self[key]
        except (KeyError, AttributeError):
            return default


class SchemaDict(DotDict):
    """
    A DotDict that ensures default attributes are always initialized
    while seamlessly accepting any extra user-defined keys.
    """
    def __init__(self, **kwargs):
        super().__init__(**kwargs)


class ComponentState(SchemaDict):
    """
    Base class for sub-states that represent collections of assets/components
    stored as parallel NumPy arrays and lists.

    Provides standard utility methods:
    - len(state): returns length based on NumPy array or list dimensions.
    - state.subset(mask): returns a sliced copy of the component state.
    - state.append(other_state): appends another component state in-place.
    """

    def __len__(self) -> int:
        for val in self.values():
            if isinstance(val, np.ndarray):
                return val.shape[0] if val.ndim >= 1 else 0
            elif isinstance(val, list):
                return len(val)
        return 0

    def subset(self, mask: Any) -> "ComponentState":
        """
        Slices all 1D/2D NumPy array fields and list fields by the given mask
        and returns a new instance of the same class.
        """
        cls = self.__class__
        new_kwargs = {}
        n_total = len(self)

        for key, val in self.items():
            if isinstance(val, np.ndarray):
                if val.ndim == 1:
                    new_kwargs[key] = val[mask]
                elif val.ndim == 2:
                    if val.shape[0] == n_total:
                        new_kwargs[key] = val[mask, :]
                    elif val.shape[1] == n_total:
                        new_kwargs[key] = val[:, mask]
                    else:
                        new_kwargs[key] = val[mask]
                else:
                    new_kwargs[key] = val[mask]
            elif isinstance(val, list):
                if isinstance(mask, np.ndarray) and mask.dtype == bool:
                    indices = np.where(mask)[0]
                elif isinstance(mask, slice):
                    indices = range(*mask.indices(len(val)))
                else:
                    indices = mask
                new_kwargs[key] = [val[i] for i in indices]
            else:
                new_kwargs[key] = val

        return cls(**new_kwargs)

    def append(self, other: "ComponentState") -> None:
        """
        Appends all matching NumPy arrays and list attributes from `other` in-place.
        """
        for key, val in other.items():
            if key in self:
                target = self[key]
                if isinstance(target, np.ndarray) and isinstance(val, np.ndarray):
                    if target.ndim == 1:
                        self[key] = np.append(target, val)
                    elif target.ndim == 2:
                        if target.shape[1] == val.shape[1]:
                            self[key] = np.vstack([target, val])
                        elif target.shape[0] == val.shape[0]:
                            self[key] = np.hstack([target, val])
                        else:
                            self[key] = np.vstack([target, val])
                elif isinstance(target, list) and isinstance(val, list):
                    target.extend(val)


class CountsState(SchemaDict):
    """
    CountsState provides access to asset counts.
    If sim_data reference is available, it dynamically returns
    the actual length of component arrays.
    """
    def __init__(self, parent_sim=None, **kwargs):
        defaults = {
            'celestial': 0,
            'satellites': 0,
            'observatories': 0,
            'fixedpoints': 0,
        }
        dict.__setattr__(self, '_parent_sim', parent_sim)
        super().__init__(**{**defaults, **kwargs})

    def __getitem__(self, item):
        if hasattr(self, '_parent_sim') and self._parent_sim is not None:
            comp = self._parent_sim.get(item)
            if comp is not None and hasattr(comp, '__len__'):
                return len(comp)
        return super().__getitem__(item)

    def __getattr__(self, name):
        if name.startswith('_'):
            return super().__getattr__(name)
        if hasattr(self, '_parent_sim') and self._parent_sim is not None:
            comp = self._parent_sim.get(name)
            if comp is not None and hasattr(comp, '__len__'):
                return len(comp)
        return super().__getattr__(name)

    def items(self):
        return [(k, self[k]) for k in self.keys() if not k.startswith('_')]

    def values(self):
        return [self[k] for k in self.keys() if not k.startswith('_')]


class CelestialState(ComponentState):
    def __init__(self, **kwargs):
        defaults = {
            'position': np.zeros((2, 3)),
            'velocity': np.zeros((2, 3)),
            'acceleration': np.zeros((2, 3)),
        }
        super().__init__(**{**defaults, **kwargs})


class SatellitesState(ComponentState):
    def __init__(self, **kwargs):
        defaults = {
            'position': np.zeros((0, 3)),
            'velocity': np.zeros((0, 3)),
            'acceleration': np.zeros((0, 3)),
            'orbital_elements': np.zeros((0, 6)),
            'epochs': [],
        }
        super().__init__(**{**defaults, **kwargs})


class FixedPointsState(ComponentState):
    def __init__(self, **kwargs):
        defaults = {
            'position': np.zeros((0, 3)),
            'exclusion': np.zeros((0,)),
            'size': np.zeros((0,)),
            'albedo': np.zeros((0,)),
        }
        super().__init__(**{**defaults, **kwargs})

    def add_target(self, position: np.ndarray, size: float, albedo: float = 0.2) -> None:
        pos = np.asarray(position, dtype=float).reshape(-1, 3)
        n_new = pos.shape[0]
        self.position = np.vstack([self.position, pos]) if self.position.size else pos
        self.exclusion = np.append(self.exclusion, np.zeros(n_new, dtype=int))
        self.size = np.append(self.size, np.full(n_new, size, dtype=float))
        self.albedo = np.append(self.albedo, np.full(n_new, albedo, dtype=float))


class ObservatoriesState(ComponentState):
    def __init__(self, **kwargs):
        defaults = {
            'position': np.zeros((0, 3)),
            'velocity': np.zeros((0, 3)),
            'acceleration': np.zeros((0, 3)),
            'pointing': np.zeros((0, 3)),
            'latitude': np.zeros(0),
            'longitude': np.zeros(0),
            'altitude': np.zeros(0),
        }
        super().__init__(**{**defaults, **kwargs})


@dataclass
class CadenceGroup:
    """
    Holds scheduling information for one group of detectors that share
    an identical integration time.

    Attributes:
        scanInterval: Integration period in seconds for this group.
        scanMask:     Boolean array (length = total detectors) selecting
                      the detectors that belong to this group.
        scanNext:     The datetime at which the next scan for this group
                      is due.  Initialised to sim_data.time so that every
                      group fires on the first call to nextIntegration.
    """
    scanInterval: float
    scanMask: np.ndarray
    scanNext: datetime


class CadenceState(SchemaDict):
    """
    Top-level cadence schedule stored in sim_data.cadenceStructure.

    Attributes:
        cadenceList: Ordered list of CadenceGroup objects, one per unique
                     integration time found in sim_data.detector.
        nextTime:    The datetime of the earliest upcoming scan across all
                     groups.
        nextGroup:   Index into cadenceList identifying which group fires
                     next.
    """
    def __init__(self, **kwargs):
        defaults = {
            'cadenceList': [],
            'nextTime': None,
            'nextGroup': 0,
        }
        super().__init__(**{**defaults, **kwargs})


class SimulationState(SchemaDict):
    """
    SimulationState models the main simulation data structure.
    It inherits from SchemaDict to maintain 100% backward compatibility
    with dict-style lookup/subscription and isinstance(..., dict) checks,
    but supports attribute-style access for better structure.
    """
    _FIELD_SCHEMAS = {
        'counts': CountsState,
        'celestial': CelestialState,
        'satellites': SatellitesState,
        'fixedpoints': FixedPointsState,
        'observatories': ObservatoriesState,
        'cadenceStructure': CadenceState,
    }

    def __init__(self, **kwargs):
        if 'start_time' not in kwargs:
            raise TypeError("SimulationState requires a 'start_time' parameter.")
        start = kwargs['start_time']

        defaults = {
            'time': start,
            'delta_time': 60.0,
            'counts': None,
            'pointing_spheres': {},
            'detector': None,
            'satellites': None,
            'observatories': None,
            'fixedpoints': None,
            'celestial': None,
            'cadenceStructure': None,
        }
        super().__init__(**{**defaults, **kwargs})
        self.counts = CountsState(parent_sim=self)

    def __setitem__(self, key, value):
        if key == 'detector':
            from detector import DetectorArray
            cls = DetectorArray
        else:
            cls = self._FIELD_SCHEMAS.get(key)

        if cls is not None and isinstance(value, dict) and not isinstance(value, cls):
            value = cls(**value)
        super().__setitem__(key, value)

    def get_detector_positions(self, mask=None) -> np.ndarray:
        """
        Builds and returns an (N, 3) NumPy array of observer/detector positions
        resolved by detector.category and detector.asset_index.

        Args:
            mask: Optional boolean or index array to subset the output.

        Returns:
            (N, 3) array of GCRS positions for detectors.
        """
        if self.detector is None or len(self.detector) == 0:
            return np.zeros((0, 3), dtype=float)

        det = self.detector
        num_detectors = len(det)
        category_array = np.array(det.category)
        asset_index_array = det.asset_index

        _pos_map = {}
        if self.satellites and len(self.satellites) > 0:
            _pos_map['satellites'] = self.satellites.position
        if self.observatories and len(self.observatories) > 0:
            _pos_map['observatories'] = self.observatories.position

        all_positions = np.zeros((num_detectors, 3), dtype=float)
        for i in range(num_detectors):
            cat = category_array[i]
            pos_array = _pos_map.get(cat)
            if pos_array is not None:
                all_positions[i] = pos_array[asset_index_array[i]]

        if mask is not None:
            return all_positions[mask]
        return all_positions


def create_empty_simulation(start_time: datetime, delta_time: float = 60.0) -> SimulationState:
    """
    Initializes a minimal, empty data structure for a space simulation.

    Args:
        start_time: The starting time and date of the simulation. This must be a
                    timezone-aware datetime object set to UTC.
        delta_time: The time step for the simulation in seconds.

    Returns:
        A SimulationState object representing the basic simulation state, not yet filled except
        for date and time and counts and pointing spheres.
    """
    if not isinstance(start_time, datetime):
        raise TypeError("start_time must be a datetime object.")
    if start_time.tzinfo is None:
        raise ValueError("start_time must be timezone-aware. Please set tzinfo.")

    return SimulationState(
        start_time=start_time,
        time=start_time,
        delta_time=delta_time
    )

\end{lstlisting}



\chapter{constants.py: Physical \& Array Constants}
\label{sec:constants}
\index{constants.py}

\section*{Overview}
Contains physical constants (Earth radius, AU, angular conversions) and element indices.

\section*{Literate Source Code}
The source code of \texttt{constants.py} is documented below. Section label: \ref{sec:constants}.

\begin{lstlisting}[language=Python,caption={constants.py Source Code},label={code:constants}]
# --- Global Constants for Array Indices ---
# These constants define column indices for numpy arrays, making the
# code more readable and preventing errors from using "magic numbers".
from math import pi


# -- Radii in Meters --
EARTH_RADIUS = 6378137.0
MOON_RADIUS = 1737400.0
GEO_RADIUS = 42164140.0
MOON_ORBIT_RADIUS = 384400000.0


# -- Some UsefulConstants --
ARCSEC = 2.0 *pi/(360*3600)
DEGREE = 3600 * ARCSEC



# -- Orbital Elements Array Indices --
ORBITAL_A_IDX = 0              # Semi-major axis in meters
ORBITAL_E_IDX = 1              # Eccentricity (dimensionless)
ORBITAL_I_IDX = 2              # Inclination in radians
ORBITAL_RAAN_IDX = 3           # Right Ascension of the Ascending Node in radians
ORBITAL_ARGP_IDX = 4           # Argument of Perigee in radians
ORBITAL_M_IDX = 5              # Mean Anomaly in radians

# -- Pointing State Array Indices --
POINTING_COUNT_IDX = 0         # Number of points in the pointing grid
POINTING_PLACE_IDX = 1         # Current index in the pointing sequence

\end{lstlisting}



\chapter{detector.py: Optoelectronic Detector Array}
\label{sec:detector}
\index{detector.py}

\section*{Overview}
Defines DetectorArray (subclass of ComponentState) and detector construction/slicing functions.

\section*{Literate Source Code}
The source code of \texttt{detector.py} is documented below. Section label: \ref{sec:detector}.

\begin{lstlisting}[language=Python,caption={detector.py Source Code},label={code:detector}]
from typing import List, Any
from constants import *
import numpy as np
import math
from minimalsimulation import SchemaDict, ComponentState
from radiometry_data import FILTER_DATA
from radiometry_calcs import *


class DetectorArray(ComponentState):
    """
    DetectorArray models detector parameters for all satellites.

    It inherits from SchemaDict so it supports both attribute-style
    (dot notation) and dictionary subscript access. All array-valued fields
    have a length equal to the number of detectors (satellites).

    Fields:
        apertureArea:    Entrance-aperture area in m^2 (1-D array).
        pixelOmega:      Pixel solid angle in sr, assuming square pixels
                         (1-D array).
        qe:              Total system quantum efficiency from entrance
                         aperture to detector (1-D array).
        photoEff:        Fraction of light captured in the photometry
                         aperture (1-D array).
        pixCount:        Number of pixels in the detector array (1-D).
        solarEx:         Solar exclusion angle in radians (1-D array).
        lunarEx:         Lunar exclusion angle in radians (1-D array).
        earthEx:         Earth limb exclusion angle in radians (1-D array).
        skyBack:         Sky background flux in photons/s/m^2/sr (1-D).
        zpCal:           Filter zero-point calibration (1-D array).
        integrationTime: Integration time in seconds (1-D array).
        fov:             Full field-of-view diameter in radians (1-D).
        ifov:            Pixel field-of-view in radians (1-D array).
        filt:            Filter band name for each detector (list of str).
        category:        The asset category this detector belongs to (e.g. 'satellites' or 'observatories') (list of str).
        asset_index:     The index of the asset within its category (1-D array of int).
        pointing:        Pointing unit vector for each detector,
                         shape (n, 3).
        pointing_state:  Pointing scheduler state, shape (2, n).
                         Row 0 = total grid points, row 1 = current index.
    """
    def __init__(self, **kwargs):
        defaults = {
            'apertureArea': np.zeros(0, dtype=float),
            'pixelOmega': np.zeros(0, dtype=float),
            'qe': np.zeros(0, dtype=float),
            'photoEff': np.zeros(0, dtype=float),
            'pixCount': np.zeros(0, dtype=float),
            'solarEx': np.zeros(0, dtype=float),
            'lunarEx': np.zeros(0, dtype=float),
            'earthEx': np.zeros(0, dtype=float),
            'skyBack': np.zeros(0, dtype=float),
            'zpCal': np.zeros(0, dtype=float),
            'integrationTime': np.zeros(0, dtype=float),
            'fov': np.zeros(0, dtype=float),
            'ifov': np.zeros(0, dtype=float),
            'filt': [],
            'category': [],
            'asset_index': np.zeros(0, dtype=int),
            'pointing': np.zeros((0, 3), dtype=float),
            'pointing_state': np.zeros((2, 0), dtype=int),
        }
        super().__init__(**{**defaults, **kwargs})


def setDetectorFOV(sim_data: Any, fovSize: float) -> None:
    """
    Sets the field-of-view of all detectors in the simulation to a single size.

    This is an ad-hoc helper function mainly used for testing, not regular operation.

    Args:
        sim_data: The main simulation data structure.
        fovSize: The new field-of-view diameter in radians.
    """
    det = sim_data.detector if hasattr(sim_data, 'detector') else sim_data['detector']
    count = len(det.fov)
    det.fov = np.full(count, fovSize)


def setDetectorIntegrationTime(sim_data: Any, itime: float) -> None:
    """
    Sets the integration time of all detectors in the simulation to a single value.

    This is an ad-hoc helper function mainly used for testing, not regular operation.

    Args:
        sim_data: The main simulation data structure.
        itime: The new integration time in seconds.
    """
    det = sim_data.detector if hasattr(sim_data, 'detector') else sim_data['detector']
    count = len(det.fov)
    det.integrationTime = np.full(count, itime)


##########################################################

def makeBlankDetector(n: int) -> DetectorArray:
    """
    makeBlankDetector makes and returns a DetectorArray
    with parameters initialized to zero arrays of length n.
    """
    return DetectorArray(
        apertureArea=np.zeros(n, dtype=float),
        pixelOmega=np.zeros(n, dtype=float),
        qe=np.zeros(n, dtype=float),
        photoEff=np.zeros(n, dtype=float),
        pixCount=np.zeros(n, dtype=float),
        solarEx=np.zeros(n, dtype=float),
        lunarEx=np.zeros(n, dtype=float),
        earthEx=np.zeros(n, dtype=float),
        skyBack=np.zeros(n, dtype=float),
        zpCal=np.zeros(n, dtype=float),
        integrationTime=np.zeros(n, dtype=float),
        fov=np.zeros(n, dtype=float),
        ifov=np.zeros(n, dtype=float),
        filt=[""] * n,
        category=[""] * n,
        asset_index=np.zeros(n, dtype=int),
        pointing=np.zeros((n, 3), dtype=float),
        pointing_state=np.zeros((2, n), dtype=int)
    )

##########################################################

def makeDetector(n, band, fov, ifov, aper, intTime: float = 1.0,
                 qe=0.5, photfrac=0.7,
                 solarex=20.0 * DEGREE, lunarex=10.0 * DEGREE,
                 earthex=15.0 * DEGREE, category=None, asset_index=None):
    '''
    makeDetector takes parameters of a sensor and stuffs a filter array
    and a detector array, which it returns.

    Args:
        n:        Number of sensors to produce.
        band:     Band the measurement takes place in (see radiometry_data).
        fov:      Full field-of-view diameter in radians.
        ifov:     Pixel field-of-view diameter in radians.
        aper:     Aperture diameter - assumed to be a disk - in meters.
                  The actual value stored and used is the apertureArea.
        intTime:  Integration time in seconds.  Defaults to 1.0.
        qe:       Quantum efficiency of the system from entrance aperture
                  to detector.
        photfrac: Fraction of light captured in the photometry aperture.
        solarex:  Solar exclusion angle in radians.
        lunarex:  Lunar exclusion angle in radians.
        earthex:  Earth limb exclusion angle in radians.

    Creates a detector for any platform (space or ground).
    '''
    detect = makeBlankDetector(n)
    detect.apertureArea[:] = math.pi * (aper / 2) ** 2
    detect.pixelOmega[:] = ifov ** 2
    detect.qe[:] = qe
    detect.photoEff[:] = photfrac
    pixels = (fov / ifov) ** 2
    detect.pixCount[:] = pixels
    detect.solarEx[:] = solarex
    detect.lunarEx[:] = lunarex
    detect.earthEx[:] = earthex
    detect.skyBack[:] = (
        amag(FILTER_DATA[band]['sky'])
        * FILTER_DATA[band]['zero_point']
        / (ARCSEC ** 2)
    )
    detect.zpCal[:] = FILTER_DATA[band]['zero_point']
    detect.integrationTime[:] = intTime
    detect.fov[:] = fov
    detect.ifov[:] = ifov
    detect.filt = [band] * n
    detect.category = category if category is not None else [""] * n
    detect.asset_index = asset_index if asset_index is not None else np.zeros(n, dtype=int)
    detect.pointing = np.zeros((n, 3), dtype=float)
    detect.pointing_state = np.zeros((2, n), dtype=int)
    return detect


def requiredIntegrationTime(limitingMag, SNR, d, debug=0):
    '''
    Calculates the required integration time to achieve a given limiting
    magnitude with a specified signal-to-noise ratio (SNR).

    Args:
        limitingMag (float): The desired limiting magnitude.
        SNR (float):         The target signal-to-noise ratio.
        d (DetectorArray):   A detector object containing the required
                             parameters.
        debug (int, optional): If set to 1, prints the intermediate
                               variables used in the calculation.
                               Defaults to 0.

    This function is based on the radiometric equation, solving for
    the integration time t:
        SNR = (Signal) / sqrt(Signal + Background)
        t = SNR^2 * (beta * omega) / (alpha^2 * A * eta * f^2)

    Returns:
        float: The required integration time in seconds.
    '''
    gamma = SNR
    beta = d.skyBack
    omega = d.pixelOmega
    alpha = amag(limitingMag) * d.zpCal
    A = d.apertureArea
    eta = d.qe
    f = d.photoEff
    if debug == 1:
        print(f"gamma {gamma:.2e}")
        print("beta is,", beta)
        print("omega", omega)
        print("alpha is", alpha)
        print('A', A)
        print('eta', eta)
        print('f', f)
    t = gamma ** 2 * beta * omega / (alpha ** 2 * A * eta * (f ** 2))
    return t


def appendDetector(cd: DetectorArray, new_cd: DetectorArray) -> None:
    """
    Appends the attributes of new_cd to the existing detector object cd
    in-place.

    Iterates over all declared DetectorArray dataclass fields.  Numpy
    1-D arrays are concatenated with np.append; 2-D arrays are joined
    with np.vstack (or np.hstack for pointing_state which is (2, n));
    list fields are extended.
    """
    for attr_name in cd.keys():
        new_val = getattr(new_cd, attr_name)
        cur_val = getattr(cd, attr_name)
        if isinstance(cur_val, np.ndarray):
            if cur_val.ndim == 1:
                setattr(cd, attr_name, np.append(cur_val, new_val))
            elif cur_val.ndim == 2:
                if attr_name == 'pointing_state':
                    setattr(cd, attr_name, np.hstack([cur_val, new_val]))
                else:
                    setattr(cd, attr_name, np.vstack([cur_val, new_val]))
        elif isinstance(cur_val, list):
            cur_val.extend(new_val)


###########################################################

def testdetector():
    '''
    testdetector creates an example that can be compared
    with some of the stuff in Curio.
    '''
    detect = makeDetector(1, "V", 1 * DEGREE,
                          2 * ARCSEC, 1,
                          qe=0.2, photfrac=1.0)
    print(requiredIntegrationTime(20.5, 10, detect, debug=1))

if __name__ == '__main__':
    testdetector()

\end{lstlisting}



\chapter{propagation.py: Satellite Orbit Propagation}
\label{sec:propagation}
\index{propagation.py}

\section*{Overview}
Implements SGP4 TLE loading and Keplerian orbit propagation for satellites.

\section*{Literate Source Code}
The source code of \texttt{propagation.py} is documented below. Section label: \ref{sec:propagation}.

\begin{lstlisting}[language=Python,caption={propagation.py Source Code},label={code:propagation}]
import numpy as np
from datetime import datetime, timezone
from typing import Dict, Any, List, Tuple
from sgp4.api import Satrec
from astropy.time import Time
from astropy.coordinates import get_body, GCRS
import astropy.units as u

from constants import *



def add_satellites_from_tle(sim_data: Any, tle_file_path: str, sat_category: str) -> None:
    """
    Adds and initializes a category of satellites from a TLE file.
    Note although the TLEs are loaded, positions etc. are not.

    Args:
        sim_data: The main simulation data structure.
        tle_file_path: Path to the TLE file.
        sat_category: The key for this satellite category (e.g., 'satellites').

    Data added to the set_category element of sim_data
    position
    velocity
    acceleration
    orbital_elements
    epochs
    """
    from detector import makeBlankDetector, appendDetector
    orbital_elements, epochs = readtle(tle_file_path)
    num_sats = len(epochs)

    sim_data.counts[sat_category] = num_sats
    new_detector = makeBlankDetector(num_sats)
    new_detector.category = [sat_category] * num_sats
    new_detector.asset_index = np.arange(num_sats, dtype=int)
    if 'detector' not in sim_data or not sim_data.detector:
        sim_data.detector = new_detector
    else:
        appendDetector(sim_data.detector, new_detector)
    sim_data[sat_category] = {
        'position': np.zeros((num_sats, 3), dtype=float),
        'velocity': np.zeros((num_sats, 3), dtype=float),
        'acceleration': np.zeros((num_sats, 3), dtype=float),
        'orbital_elements': orbital_elements,
        'epochs': epochs,
    }

def readtle(tle_file_path: str) -> Tuple[np.ndarray, List[datetime]]:
    """
    Reads a TLE file and extracts orbital elements and epochs for each satellite.

    The array returned has the orbital elements in "canonical" order.
    """
    orbital_elements_list = []
    epochs_list = []
    with open(tle_file_path, 'r') as f:
        lines = f.readlines()

    for i in range(0, len(lines), 3):
        line1 = lines[i+1].strip()
        line2 = lines[i+2].strip()

        satellite = Satrec.twoline2rv(line1, line2)

        jd, fr = satellite.jdsatepoch, satellite.jdsatepochF
        epoch_dt = Time(jd, fr, format='jd', scale='utc').to_datetime(timezone.utc)
        epochs_list.append(epoch_dt)

        a = satellite.a * satellite.radiusearthkm * 1000.0
        e = satellite.ecco
        inc = satellite.inclo
        raan = satellite.nodeo
        argp = satellite.argpo
        M = satellite.mo

        elements = np.zeros(6)
        elements[ORBITAL_A_IDX] = a
        elements[ORBITAL_E_IDX] = e
        elements[ORBITAL_I_IDX] = inc
        elements[ORBITAL_RAAN_IDX] = raan
        elements[ORBITAL_ARGP_IDX] = argp
        elements[ORBITAL_M_IDX] = M
        orbital_elements_list.append(elements)

    return np.array(orbital_elements_list, dtype=float), epochs_list

def propagate_satellites(data_struct: Any, time_date: datetime, sat_category: str = None) -> Any:
    """
    data_struct is the "standard" structure for the simulation
    time_date - a datetime - is the time that the satellites are propagated to.
    NOTE that the datetime structure in simulation_data is NOT updated to this value by this function!

    Updates satellite positions and pointing vectors based on their orbital elements.
    
    """
    MU_EARTH = 3.986004418e14
    time_date_timestamp = time_date.timestamp()

    if sat_category:
        categories = [sat_category]
    elif hasattr(data_struct, 'counts') and data_struct.counts:
        non_sat_keys = {'celestial', 'observatories', 'fixedpoints'}
        categories = [
            k for k, v in data_struct.counts.items()
            if k not in non_sat_keys and v > 0
        ]
        if not categories:
            categories = ['satellites']
    else:
        categories = ['satellites']

    for category in categories:
        if category not in data_struct.counts or data_struct.counts[category] == 0:
            continue

        elements = data_struct[category].orbital_elements
        epochs = data_struct[category].epochs

        epoch_timestamps = np.array([e.timestamp() for e in epochs])
        delta_t_array = time_date_timestamp - epoch_timestamps

        a = elements[:, ORBITAL_A_IDX]
        e = elements[:, ORBITAL_E_IDX]
        i = elements[:, ORBITAL_I_IDX]
        raan = elements[:, ORBITAL_RAAN_IDX]
        argp = elements[:, ORBITAL_ARGP_IDX]
        M0 = elements[:, ORBITAL_M_IDX]

        # Guard against uninitialized/invalid orbits (a <= 0 or e >= 1.0)
        safe_a = np.where(a <= 0, EARTH_RADIUS + 400000.0, a)
        safe_e = np.clip(e, 0.0, 0.999999)

        n = np.sqrt(MU_EARTH / safe_a**3)
        M = (M0 + n * delta_t_array) % (2 * np.pi)

        E = M.copy()
        for _ in range(10):
            f_E = E - safe_e * np.sin(E) - M
            f_prime_E = 1 - safe_e * np.cos(E)
            f_prime_E[f_prime_E == 0] = 1e-10
            E = E - f_E / f_prime_E

        tan_nu_half = np.sqrt((1 + safe_e) / (1 - safe_e)) * np.tan(E / 2)
        nu = 2 * np.arctan(tan_nu_half)

        r = safe_a * (1 - safe_e * np.cos(E))

        x_pqw = r * np.cos(nu)
        y_pqw = r * np.sin(nu)

        cos_raan = np.cos(raan)
        sin_raan = np.sin(raan)
        cos_argp = np.cos(argp)
        sin_argp = np.sin(argp)
        cos_i = np.cos(i)
        sin_i = np.sin(i)

        P_x = cos_argp * cos_raan - sin_argp * sin_raan * cos_i
        P_y = cos_argp * sin_raan + sin_argp * cos_raan * cos_i
        P_z = sin_argp * sin_i

        Q_x = -sin_argp * cos_raan - cos_argp * sin_raan * cos_i
        Q_y = -sin_argp * sin_raan + cos_argp * cos_raan * cos_i
        Q_z = cos_argp * sin_i

        x_gcrs = x_pqw * P_x + y_pqw * Q_x
        y_gcrs = x_pqw * P_y + y_pqw * Q_y
        z_gcrs = x_pqw * P_z + y_pqw * Q_z

        positions = np.vstack((x_gcrs, y_gcrs, z_gcrs)).T
        data_struct[category].position = positions

        norms = np.linalg.norm(positions, axis=1)[:, np.newaxis]
        norms[norms == 0] = 1.0


    return data_struct

\end{lstlisting}



\chapter{observatories.py: Ground Station Observatories}
\label{sec:observatories}
\index{observatories.py}

\section*{Overview}
Implements GCRS coordinate conversion and sidereal velocity calculations for ground stations.

\section*{Literate Source Code}
The source code of \texttt{observatories.py} is documented below. Section label: \ref{sec:observatories}.

\begin{lstlisting}[language=Python,caption={observatories.py Source Code},label={code:observatories}]
import numpy as np
from typing import Dict, Any

def add_observatories(
    sim_data: Any,
    num_observatories: int,
    latitudes: np.ndarray = None,
    longitudes: np.ndarray = None,
    altitudes: np.ndarray = None
) -> None:
    """
    Adds observatory data structures to the simulation data.

    Args:
        sim_data: The main simulation data structure.
        num_observatories: The number of observatories to add.
        latitudes: Latitude coordinates in degrees.
        longitudes: Longitude coordinates in degrees.
        altitudes: Altitude coordinates in meters.
    """
    if not isinstance(num_observatories, int) or num_observatories < 0:
        raise ValueError("num_observatories must be a non-negative integer.")

    if latitudes is None:
        latitudes = np.zeros(num_observatories)
    if longitudes is None:
        longitudes = np.zeros(num_observatories)
    if altitudes is None:
        altitudes = np.zeros(num_observatories)

    from detector import makeBlankDetector, appendDetector
    detector = makeBlankDetector(num_observatories)
    detector.category = ['observatories'] * num_observatories
    
    from minimalsimulation import ObservatoriesState
    if sim_data.counts.get('observatories', 0) > 0:
        # Append to existing
        existing = sim_data.observatories
        detector.asset_index = np.arange(sim_data.counts.observatories, sim_data.counts.observatories + num_observatories, dtype=int)
        sim_data.counts.observatories += num_observatories
        
        existing.latitude = np.append(existing.latitude, np.array(latitudes, dtype=float))
        existing.longitude = np.append(existing.longitude, np.array(longitudes, dtype=float))
        existing.altitude = np.append(existing.altitude, np.array(altitudes, dtype=float))
        existing.position = np.vstack([existing.position, np.zeros((num_observatories, 3), dtype=float)])
        existing.velocity = np.vstack([existing.velocity, np.zeros((num_observatories, 3), dtype=float)])
        existing.acceleration = np.vstack([existing.acceleration, np.zeros((num_observatories, 3), dtype=float)])
        existing.pointing = np.vstack([existing.pointing, np.zeros((num_observatories, 3), dtype=float)])
    else:
        sim_data.counts.observatories = num_observatories
        detector.asset_index = np.arange(num_observatories, dtype=int)
        sim_data.observatories = ObservatoriesState(
            latitude=np.array(latitudes, dtype=float),
            longitude=np.array(longitudes, dtype=float),
            altitude=np.array(altitudes, dtype=float),
            position=np.zeros((num_observatories, 3), dtype=float),
            velocity=np.zeros((num_observatories, 3), dtype=float),
            acceleration=np.zeros((num_observatories, 3), dtype=float),
            pointing=np.zeros((num_observatories, 3), dtype=float),
        )
    
    if 'detector' not in sim_data or not sim_data.detector:
        sim_data.detector = detector
    else:
        appendDetector(sim_data.detector, detector)

    # Initial propagation of observatory positions to the current simulation time
    propagate_observatories(sim_data, sim_data.time)


def propagate_observatories(sim_data: Any, time_date: Any) -> None:
    """
    Propagates observatory positions in GCRS (ECI) based on Earth rotation.
    """
    obs = sim_data.observatories
    if not obs or len(obs.latitude) == 0:
        return

    from astropy.coordinates import EarthLocation
    from astropy.time import Time
    import astropy.units as u

    # Define locations
    locations = EarthLocation(
        lat=obs.latitude * u.deg,
        lon=obs.longitude * u.deg,
        height=obs.altitude * u.m
    )

    t = Time(time_date)
    # Get coordinates in Geocentric Celestial Reference System (GCRS)
    gcrs_coords = locations.get_gcrs(t)

    # GCRS coordinates are in meters
    x = gcrs_coords.cartesian.x.to(u.m).value
    y = gcrs_coords.cartesian.y.to(u.m).value
    z = gcrs_coords.cartesian.z.to(u.m).value

    # If x, y, z are scalars (e.g. single observatory), stack them correctly
    if np.isscalar(x):
        obs.position = np.array([[x, y, z]], dtype=float)
    else:
        obs.position = np.vstack([x, y, z]).T

    # Compute velocity via finite-difference of GCRS positions.
    # The naive omegaxr formula is only valid in the ECEF/ITRS
    # (body-fixed) frame, but obs.position is in GCRS (inertial).
    # A small Deltat finite-difference gives the correct GCRS velocity
    # including precession and nutation effects.
    from datetime import timedelta
    dt_seconds = 1.0
    t2 = Time(time_date + timedelta(seconds=dt_seconds))
    gcrs_coords2 = locations.get_gcrs(t2)

    x2 = gcrs_coords2.cartesian.x.to(u.m).value
    y2 = gcrs_coords2.cartesian.y.to(u.m).value
    z2 = gcrs_coords2.cartesian.z.to(u.m).value

    if np.isscalar(x2):
        pos2 = np.array([[x2, y2, z2]], dtype=float)
    else:
        pos2 = np.vstack([x2, y2, z2]).T

    obs.velocity = (pos2 - obs.position) / dt_seconds

\end{lstlisting}



\chapter{celestialbodies.py: Celestial Bodies Ephemeris}
\label{sec:celestialbodies}
\index{celestialbodies.py}

\section*{Overview}
Updates Sun and Moon positions in GCRS coordinates over simulation time.

\section*{Literate Source Code}
The source code of \texttt{celestialbodies.py} is documented below. Section label: \ref{sec:celestialbodies}.

\begin{lstlisting}[language=Python,caption={celestialbodies.py Source Code},label={code:celestialbodies}]
import numpy as np
from typing import Dict, Any, Optional
from datetime import datetime
from astropy.time import Time
from astropy.coordinates import get_body, GCRS
import astropy.units as u

def add_celestial_bodies(sim_data: Any) -> None:
    """
    Adds celestial body structures (for Sun and Moon) to the simulation data.

    Args:
        sim_data: The simulation data structure.
    """
    sim_data.counts.celestial = 2  # Sun and Moon
    sim_data.celestial = {
        'position': np.zeros((2, 3), dtype=float),
        'velocity': np.zeros((2, 3), dtype=float),
        'acceleration': np.zeros((2, 3), dtype=float),
    }

def celestial_update(data_struct: Any, time_date: Optional[datetime] = None) -> Any:
    """
    Calculates and updates the positions of the Sun and Moon.
    """
    if time_date is None:
        time_date = data_struct.time

    if time_date.tzinfo is None:
        raise ValueError("time_date must be timezone-aware.")

    astro_time = Time(time_date)
    sun_coords = get_body("sun", astro_time)
    sun_gcrs = sun_coords.transform_to(GCRS(obstime=astro_time))
    moon_coords = get_body("moon", astro_time)
    moon_gcrs = moon_coords.transform_to(GCRS(obstime=astro_time))

    celestial_pos = data_struct.celestial.position
    celestial_pos[0] = sun_gcrs.cartesian.xyz.to(u.m).value
    celestial_pos[1] = moon_gcrs.cartesian.xyz.to(u.m).value

    return data_struct

\end{lstlisting}



\chapter{targets.py: Reference Target Generation}
\label{sec:targets}
\index{targets.py}

\section*{Overview}
Generates logarithmic spherical distributions of fixed reference targets.

\section*{Literate Source Code}
The source code of \texttt{targets.py} is documented below. Section label: \ref{sec:targets}.

\begin{lstlisting}[language=Python,caption={targets.py Source Code},label={code:targets}]
import numpy as np
# from datetime import datetime, timezone
from typing import Dict, Any
from generate_log_spherical_points import generate_log_spherical_points
from constants import *
from minimalsimulation import FixedPointsState

def add_fixed_points(
    sim_data: Any,
    num_points: int = 100,
    size: float = 1.0,
    innerRadius: float = 2000000,
    outerRadius: float = 2 * GEO_RADIUS
) -> None:
    """
    Adds a structure for fixed reference points in the GCRS frame.

    Args:
        sim_data: The simulation data structure.
        num_points: The number of fixed points to generate.
        size: the size of the objects, 
        innerRadius: The minimum radius at which test points will be created.
        outerRadius: The maximum radius at which test points will be created.
    """
    positions = generate_log_spherical_points(
        num_points=num_points,
        inner_radius=innerRadius,
        outer_radius=outerRadius
    )
    if not sim_data.fixedpoints:
        sim_data.fixedpoints = FixedPointsState()
    sim_data.fixedpoints.add_target(positions, size=size, albedo=0.2)

\end{lstlisting}



\chapter{exclusion.py: Celestial Obstruction \& Exclusion Zones}
\label{sec:exclusion}
\index{exclusion.py}

\section*{Overview}
Computes solar, lunar, and Earth limb/horizon exclusion zones for space and ground sensors.

\section*{Literate Source Code}
The source code of \texttt{exclusion.py} is documented below. Section label: \ref{sec:exclusion}.

\begin{lstlisting}[language=Python,caption={exclusion.py Source Code},label={code:exclusion}]
import numpy as np
from typing import Dict, Any, Tuple, Optional

from constants import (
    EARTH_RADIUS, MOON_RADIUS
)
def exclusion(
    data_struct: Dict[str, Any],
    satellite_index: int,
    sat_category: str = 'satellites',
    print_debug: bool = False
) -> int:
    """
    Determines if a satellite's pointing vector is excluded by the Sun, Moon, or Earth.

    This function is vectorized to check for exclusions from all three bodies (Sun,
    Moon, Earth) simultaneously for a single satellite.

    Args:
        data_struct: The main simulation data dictionary.
        satellite_index: The index of the satellite to check.
        sat_category: The satellite category (e.g. 'satellites', 'red_satellites').
        print_debug: If True, prints detailed debug information for the calculation.

    Returns:
        1 if the satellite's view is excluded by any of the bodies, 0 otherwise.
    """
    # Determine asset category and index using the detector's internal tracking
    detector_props = data_struct.detector
    asset_category = detector_props.category[satellite_index]
    asset_idx = detector_props.asset_index[satellite_index]

    sat_pos = getattr(data_struct, asset_category).position[asset_idx]
    sat_pointing = data_struct.detector.pointing[satellite_index]
    
    norm_pointing = np.linalg.norm(sat_pointing)
    if norm_pointing < 1e-9:
        return 1  # Not pointing anywhere, treat as excluded

    u_sat_pointing = sat_pointing / norm_pointing

    detector_props = data_struct.detector

    # Local horizon exclusion check for ground-based observatories
    if asset_category == 'observatories':
        norm_pos = np.linalg.norm(sat_pos)
        if norm_pos > 1e-9:
            zenith_normal = sat_pos / norm_pos
            cos_zenith = np.clip(np.dot(zenith_normal, u_sat_pointing), -1.0, 1.0)
            zenith_angle = np.arccos(cos_zenith)
            # Minimum elevation angle is stored in earthEx
            min_elevation = detector_props.earthEx[satellite_index]
            max_zenith = np.pi / 2.0 - min_elevation
            if zenith_angle > max_zenith:
                if print_debug:
                    print(f"--- Horizon Exclusion for Observatory {asset_idx} ---")
                    print(f"  Zenith Angle: {np.rad2deg(zenith_angle):.2f} deg, "
                          f"Max Zenith: {np.rad2deg(max_zenith):.2f} deg (elev < {np.rad2deg(min_elevation):.2f} deg)")
                return 1

    # Celestial body positions and radii
    body_positions = np.array([
        data_struct.celestial.position[0],  # Sun
        data_struct.celestial.position[1],  # Moon
        [0.0, 0.0, 0.0]  # Earth's center
    ])
    body_radii = np.array([0.0, MOON_RADIUS, EARTH_RADIUS])

    # Satellite-specific exclusion angles
    exclusion_angles = np.array([
        detector_props.solarEx[satellite_index],
        detector_props.lunarEx[satellite_index],
        detector_props.earthEx[satellite_index]
    ])

    # --- Vectorized Calculations ---
    vecs_to_bodies = body_positions - sat_pos
    dists_to_bodies = np.linalg.norm(vecs_to_bodies, axis=1)

    # Calculate angles to bodies, handling potential division by zero
    angles = np.full(3, np.pi)
    valid_dists = dists_to_bodies > 1e-9
    
    u_vecs_to_bodies = np.zeros_like(vecs_to_bodies)
    u_vecs_to_bodies[valid_dists] = vecs_to_bodies[valid_dists] / dists_to_bodies[valid_dists, np.newaxis]

    cos_angles = np.clip(np.dot(u_vecs_to_bodies, u_sat_pointing), -1.0, 1.0)
    angles = np.arccos(cos_angles)

    # Calculate apparent radii (for Moon and Earth)
    apparent_radii = np.zeros(3)
    non_zero_dists = (dists_to_bodies[1:] > 0)
    apparent_radii[1:][non_zero_dists] = np.arctan(body_radii[1:][non_zero_dists] / dists_to_bodies[1:][non_zero_dists])

    # Check for exclusion
    is_excluded = (angles - apparent_radii) < exclusion_angles

    if asset_category == 'observatories':
        # For ground stations, Earth limb checks are replaced by the horizon check above.
        is_excluded[2] = False

    if print_debug:
        body_names = ["Sun", "Moon", "Earth"]
        asset_name = "Observatory" if asset_category == "observatories" else "Satellite"
        print(f"--- Exclusion Debug for {asset_name} {asset_idx} ---")
        for i in range(3):
            print(f"  - {body_names[i]:<5} Flag: {is_excluded[i]}, "
                  f"Angle: {np.rad2deg(angles[i]):.2f} deg, "
                  f"Excl: {np.rad2deg(exclusion_angles[i]):.2f} deg")

    return 1 if np.any(is_excluded) else 0

def update_exclusion_table(
    data_struct: Dict[str, Any],
    print_debug_for_sat: Optional[int] = None
) -> None:
    """
    Updates the exclusion table for all detectors against all
    fixed points.  Uses the detector category / asset_index
    mapping so that satellites and observatories are handled
    uniformly.

    For space-based detectors the standard Sun / Moon / Earth-limb
    exclusion checks are applied.  For ground-based observatories
    the Earth-limb check is replaced by a local-horizon elevation
    check (the minimum elevation is stored in detector.earthEx).

    The resulting matrix is stored in
    data_struct.fixedpoints.exclusion with shape
    (num_fixed_points, num_detectors).
    """
    num_fixed_points = data_struct.counts.get('fixedpoints', 0)
    num_detectors = len(data_struct.detector.filt)

    if num_detectors == 0 or num_fixed_points == 0:
        return

    targets = data_struct.fixedpoints.position
    detector_props = data_struct.detector

    # --- Build observer positions using get_detector_positions ---
    category_array = np.array(detector_props.category)
    observer_pos = data_struct.get_detector_positions()

    # --- Pointing vectors: (num_targets, num_detectors, 3) ---
    pointing_vectors = (
        targets[:, np.newaxis, :]
        - observer_pos[np.newaxis, :, :]
    )

    norm_pointing = np.linalg.norm(pointing_vectors, axis=2)
    safe_norm = np.where(norm_pointing == 0, 1.0, norm_pointing)
    u_pointing = pointing_vectors / safe_norm[:, :, np.newaxis]

    # --- Celestial body geometry ---
    body_positions = np.array([
        data_struct.celestial.position[0],
        data_struct.celestial.position[1],
        [0.0, 0.0, 0.0]
    ])
    body_radii = np.array([0.0, MOON_RADIUS, EARTH_RADIUS])

    # vecs_to_bodies: (num_detectors, 3_bodies, 3_xyz)
    vecs_to_bodies = (
        body_positions[np.newaxis, :, :]
        - observer_pos[:, np.newaxis, :]
    )
    dists_to_bodies = np.linalg.norm(vecs_to_bodies, axis=2)
    safe_dists = np.where(
        dists_to_bodies == 0, 1.0, dists_to_bodies
    )
    u_vecs_to_bodies = (
        vecs_to_bodies / safe_dists[:, :, np.newaxis]
    )

    # Angles: (num_targets, num_detectors, 3_bodies)
    #   t = target, d = detector, b = body, i = xyz
    cos_angles = np.einsum(
        'tdi,dbi->tdb', u_pointing, u_vecs_to_bodies
    )
    cos_angles = np.clip(cos_angles, -1.0, 1.0)
    angles = np.arccos(cos_angles)

    # Apparent angular radii: (num_detectors, 3_bodies)
    apparent_radii = np.zeros((num_detectors, 3))
    apparent_radii[:, 1:] = np.arctan(
        body_radii[1:][np.newaxis, :] / safe_dists[:, 1:]
    )

    # Exclusion angles: (num_detectors, 3_bodies)
    exclusion_angles = np.vstack([
        detector_props.solarEx,
        detector_props.lunarEx,
        detector_props.earthEx
    ]).T

    # Exclusion test: (num_targets, num_detectors, 3_bodies)
    is_excluded = (
        (angles - apparent_radii[np.newaxis, :, :])
        < exclusion_angles[np.newaxis, :, :]
    )

    # --- Observatory horizon check (replaces Earth-limb) ---
    obs_mask = (category_array == 'observatories')
    if np.any(obs_mask):
        obs_indices = np.where(obs_mask)[0]
        obs_positions = observer_pos[obs_indices]
        obs_norms = np.linalg.norm(obs_positions, axis=1)
        safe_obs_norms = np.where(
            obs_norms == 0, 1.0, obs_norms
        )
        zenith_vecs = (
            obs_positions / safe_obs_norms[:, np.newaxis]
        )

        # cos(zenith angle) for every (target, obs) pair
        cos_zenith = np.einsum(
            'toi,oi->to',
            u_pointing[:, obs_indices, :],
            zenith_vecs
        )
        cos_zenith = np.clip(cos_zenith, -1.0, 1.0)
        zenith_angles = np.arccos(cos_zenith)

        min_elev = exclusion_angles[obs_indices, 2]
        max_zenith = np.pi / 2.0 - min_elev

        # Replace Earth-body exclusion with horizon check
        is_excluded[:, obs_indices, 2] = (
            zenith_angles > max_zenith[np.newaxis, :]
        )

    # Collapse body axis -> scalar excluded flag
    exclusion_matrix = np.any(is_excluded, axis=2).astype(int)

    # Zero-norm pointing -> treat as excluded
    exclusion_matrix[norm_pointing < 1e-9] = 1

    # --- Debug output ---
    if (print_debug_for_sat is not None
            and 0 <= print_debug_for_sat < num_detectors):
        original_pointing = data_struct.detector.pointing.copy()
        for j in range(num_fixed_points):
            data_struct.detector.pointing[print_debug_for_sat] = (
                pointing_vectors[j, print_debug_for_sat]
            )
            exclusion(
                data_struct,
                print_debug_for_sat,
                print_debug=True,
            )
        data_struct.detector.pointing = original_pointing

    data_struct.fixedpoints.exclusion = exclusion_matrix


\end{lstlisting}



\chapter{pointing.py: Sensor Pointing \& Re-slewing}
\label{sec:pointing}
\index{pointing.py}

\section*{Overview}
Generates spherical pointing grids, updates boresight vectors, and executes random perturbation (jerk).

\section*{Literate Source Code}
The source code of \texttt{pointing.py} is documented below. Section label: \ref{sec:pointing}.

\begin{lstlisting}[language=Python,caption={pointing.py Source Code},label={code:pointing}]
import numpy as np
from typing import Dict, Any
import math
from datetime import datetime, timezone
import plotly.graph_objects as go
import os

from minimalsimulation import create_empty_simulation
from celestialbodies import add_celestial_bodies, celestial_update
from propagation import add_satellites_from_tle, propagate_satellites
from plotting_3d import plot_3d_scatter
from constants import POINTING_COUNT_IDX, POINTING_PLACE_IDX
from fibonacciSearch import pointing_vectors, resort_vectors_by_proximity
from exclusion import exclusion


def generate_pointing_sphere(sim_data: Any, n_points: int, debug: bool = False) -> None:
    """
    Generates a pointing sphere with n_points and stores it in the sim_data.pointing_spheres[n_points]
    A pointing sphere is a 3 by n_points numpy array with the 3 representing unit vectors to be
    pointed to
    These positions will be used by the update_detector_pointing to
    point sensors incrementaly.
    The index is sim_data.pointing_spheres[n] 
    If a sphere with the same number of points already exists, this function does nothing.
    The current version of the code resorts the vector by proximity to make the sky search
    more efficient, although this is not deeply optimized yet.
    """
    if n_points not in sim_data.pointing_spheres:
        print(f"Generating pointing sphere with {n_points} points...")
        sim_data.pointing_spheres[n_points] = \
           resort_vectors_by_proximity(pointing_vectors(n_points))

    if debug:
        print(f"\n--- Debugging Pointing Sphere (n_points={n_points}) ---")
        generated_vectors = sim_data.pointing_spheres[n_points]
        print(f"Total number of pointing vectors generated: {len(generated_vectors)}")

        num_to_show = min(5, len(generated_vectors))
        print(f"Showing first {num_to_show} pointing vectors and their norms:")
        print(f"{'Index':<7} {'Vector (x, y, z)':<35} {'Norm':<10}")
        print(f"{'-------':<7} {'-----------------------------------':<35} {'----------':<10}")
        for i in range(num_to_show):
            vec = generated_vectors[i]
            norm = np.linalg.norm(vec)
            print(f"{i:<7} {str(vec):<35} {norm:<10.6f}")
        print("-------------------------------------------\n")


def update_detector_pointing(sim_data: Any, sat_category: str = 'satellites', debug: bool = False) -> None:
    """
    Updates the pointing vector for each detector, skipping excluded pointing directions.
    """
    num_detectors = len(sim_data.detector.filt)  
    if num_detectors == 0:
        return

# Bring in the appropriate pieces of the data structure for easier reference
    pointing_state = sim_data.detector.pointing_state 
#    print('pointing_state in pointing.py ', pointing_state)
    pointing_vectors_all = sim_data.detector.pointing

# Iterate over satellites
    for i in range(num_detectors):
# Place a grid of vectors to use in grid
        count = int(pointing_state[POINTING_COUNT_IDX, i])
        if count == 0:
            continue
        grid = sim_data.pointing_spheres[count]
        
        place = int(pointing_state[POINTING_PLACE_IDX, i])
        start_place = place

        
        while True:
# Move to the next place, wrap around if at end
            place += 1
            if place >= count:
                place = 0
            pointing_vectors_all[i] = grid[place]
            
            excluded = exclusion(sim_data, i, sat_category=sat_category)
            if debug:
                print(f"Detector {i}: Pointing location {place}, Excluded: {excluded != 0}")
                print(grid[place])

            if excluded == 0:
                pointing_state[POINTING_PLACE_IDX, i] = place
                break
            
            if place == start_place:
                print(f"Warning: Satellite {i} has all pointing vectors excluded.")
                pointing_state[POINTING_PLACE_IDX, i] = place
                break


def detectorPointingInitialize(sim_data: Any, grid_points: int) -> None:
    """
    Assumes that sim_data.detector is loaded, but the
    pointing part of detectors is currently empty.
    Initializes pointing and pointing_state inside detector, and
    adds a pointing sphere to sim_data.
    """
    sensorCount = len(sim_data.detector.filt)
    generate_pointing_sphere(sim_data, grid_points)
    detect = sim_data.detector
    detect.pointing_state = np.zeros((2, sensorCount), dtype=int)
    detect.pointing_state[POINTING_COUNT_IDX, :] = grid_points
    detect.pointing_state[POINTING_PLACE_IDX, :] = np.random.randint(0, grid_points - 1, size=sensorCount)
    update_detector_pointing(sim_data)


def demo_exclusion_pointing():
    """
    Demonstrates satellite pointing with a solar exclusion angle and a detector
    field of view, plotting the pointing history on a sphere.
    """
    start_time = datetime(2025, 1, 1, 12, 0, 0, tzinfo=timezone.utc)
    sim_data = create_empty_simulation(start_time)
    add_celestial_bodies(sim_data)

    # Create a temporary TLE file with a single satellite for this demo
    single_tle_content = """GEO-SAT
1 99999U 99999A   25001.00000000  .00000000  00000-0  00000-0 0  9999
2 99999   0.0001 254.9961 0000001  98.4322 261.6813  1.00271900    10"""
    
    temp_tle_path = "temp_single_sat.tle"
    with open(temp_tle_path, "w") as f:
        f.write(single_tle_content)

    add_satellites_from_tle(sim_data, temp_tle_path, 'satellites')
    
    # Clean up the temporary TLE file
    os.remove(temp_tle_path)

    # Set solar exclusion angle and detector FOV for the single satellite
    sim_data.detector.solarEx[0] = math.pi/2
    sim_data.detector.fov[0] = math.pi/5

    # Generate 400 pointing points using the module's generate_pointing_sphere
    n_points_sphere = 400
    generate_pointing_sphere(sim_data, n_points_sphere)

    # Initialize pointing_state for the single detector
    # Assuming the first detector (index 0)
    sim_data.detector.pointing_state[POINTING_COUNT_IDX, 0] = n_points_sphere
    sim_data.detector.pointing_state[POINTING_PLACE_IDX, 0] = 0 # Start at the first point

    pointed_directions_history = []
    
    # Propagate the satellite once to get an initial position
    sim_data = celestial_update(sim_data, start_time)
    sim_data = propagate_satellites(sim_data, start_time, 'satellites')
    
    initial_sat_pos = sim_data.satellites.position[0]
    
    # --- Plot initialization with unit sphere ---
    fig = go.Figure()
    # Add a sphere at the origin to represent the pointing sphere itself for context
    sphere_radius = 1.0 # Unit sphere
    u_sphere = np.linspace(0, 2 * np.pi, 50)
    v_sphere = np.linspace(0, np.pi, 50)
    x_sphere = sphere_radius * np.outer(np.cos(u_sphere), np.sin(v_sphere))
    y_sphere = sphere_radius * np.outer(np.sin(u_sphere), np.sin(v_sphere))
    z_sphere = sphere_radius * np.outer(np.ones(np.size(u_sphere)), np.cos(v_sphere))
    fig.add_trace(go.Surface(
        x=x_sphere, y=y_sphere, z=z_sphere,
        colorscale='Blues', showscale=False, opacity=0.1, name='Unit Sphere'
    ))
    
    # Add initial satellite position as a marker
    # The plot_3d_scatter function would normally plot initial_sat_pos scaled,
    # but here we are just adding it as a marker on the unit sphere for context
    # as the vectors are also on the unit sphere.
    initial_sat_direction = initial_sat_pos / np.linalg.norm(initial_sat_pos)
    fig.add_trace(go.Scatter3d(
        x=[initial_sat_direction[0]], y=[initial_sat_direction[1]], z=[initial_sat_direction[2]],
        mode='markers',
        marker=dict(size=5, color='blue'),
        name='Initial Satellite Pointing'
    ))

    for i in range(400):
        update_detector_pointing(sim_data, debug=False) # Turn off debug
        current_pointed_direction = sim_data.detector.pointing[0]
        snapshot = current_pointed_direction.copy()
#        print(f"Current pointed direction: {snapshot}") # Print current direction
        pointed_directions_history.append(snapshot)

    pointed_directions_history = np.array(pointed_directions_history)
#   print("Array pointedDirectionsHistory", pointed_directions_history)

    # Add the pointing history as markers
    colors = np.arange(len(pointed_directions_history))
    fig.add_trace(go.Scatter3d(
        x=pointed_directions_history[:, 0],
        y=pointed_directions_history[:, 1],
        z=pointed_directions_history[:, 2],
        mode='markers',
        marker=dict(
            size=3,
            color=colors,
            colorscale='Plasma', # Using a gradient color
            opacity=0.7
        ),
        name='Pointing History'
    ))

    # Connect the pointing points with lines
    fig.add_trace(go.Scatter3d(
        x=pointed_directions_history[:, 0],
        y=pointed_directions_history[:, 1],
        z=pointed_directions_history[:, 2],
        mode='lines',
        line=dict(color='grey', width=1),
        name='Pointing Path'
    ))
    
    fig.update_layout(
        title="Detector Pointing with Exclusion",
        scene=dict(
            xaxis_title='X (Unit Vector)', yaxis_title='Y (Unit Vector)', zaxis_title='Z (Unit Vector)',
            aspectmode='data'
        ),
        margin=dict(r=20, b=10, l=10, t=40)
    )

#    print(pointed_directions_history[:100])
    fig.add_annotation(
        text="Generated by: demo_exclusion_pointing",
        xref="paper", yref="paper",
        x=0.5, y=1.01,
        showarrow=False,
        font=dict(size=10, color="gray"),
        xanchor="center", yanchor="bottom"
    )
    return fig

def jerk(sim_data: Any, satellite_indices: np.ndarray) -> Any:
    """
    Moves the pointing vector of specific satellites by 0.3 radians in a
    random direction.

    This function applies a random rotation to the satellites' pointing vectors.

    Args:
        sim_data: The main simulation data structure.
        satellite_indices: The indices of the satellites to modify.

    Returns:
        The modified sim_data with the updated pointing vectors.
    """
    if satellite_indices.size == 0:
        return sim_data

    p = sim_data.detector.pointing[satellite_indices]
    p_norm = p / np.linalg.norm(p, axis=1)[:, np.newaxis]

    # Generate a random vector not parallel to p_norm
    r = np.random.randn(*p.shape)
    r -= np.sum(r * p_norm, axis=1)[:, np.newaxis] * p_norm
    k_hat = r / np.linalg.norm(r, axis=1)[:, np.newaxis]

    theta = 0.3
    cos_theta = np.cos(theta)
    sin_theta = np.sin(theta)

    # Rodrigues' rotation formula
    p_new = p_norm * cos_theta + np.cross(k_hat, p_norm) * sin_theta

    sim_data.detector.pointing[satellite_indices] = p_new

    return sim_data

\end{lstlisting}



\chapter{scandetectors.py: Target Detection \& Photometric Engine}
\label{sec:scandetectors}
\index{scandetectors.py}

\section*{Overview}
Vectorized detection engine calculating signal, background noise, and SNR for visible targets.

\section*{Literate Source Code}
The source code of \texttt{scandetectors.py} is documented below. Section label: \ref{sec:scandetectors}.

\begin{lstlisting}[language=Python,caption={scandetectors.py Source Code},label={code:scandetectors}]

import numpy as np
from lambertian import lambertiansphere, includedAngle
import radiometry_calcs
from constants import EARTH_RADIUS


def get_spherical_coords(arr: np.ndarray) -> tuple[np.ndarray, np.ndarray]:
    """
    Converts 3D Cartesian coordinates to spherical angles (theta, phi).

    Args:
        arr: An (N, 3) NumPy array of Cartesian coordinates.

    Returns:
        tuple[np.ndarray, np.ndarray]:
            - theta: The inclination angle from the z-axis in radians [0, pi].
            - phi: The azimuthal angle in the x-y plane in radians [-pi, pi].
    """
    x, y, z = arr[:, 0], arr[:, 1], arr[:, 2]
    r = np.linalg.norm(arr, axis=1)
    # theta: angle from the z-axis [0, pi]
    # We use np.divide to handle potential division by zero if r is 0
    theta = np.arccos(np.clip(z / r, -1.0, 1.0))
    # phi: angle in the x-y plane [-pi, pi]
    phi = np.arctan2(y, x)
    return theta, phi

#######################################################


def scandetectors(sim_data: dict, print_output: int = 0, mask: np.ndarray = None):
    """
    Scans for and processes detector data within the simulation.

    This function orchestrates the process of determining which fixed targets
    are visible to each satellite's detector and calculates the resulting
    signal-to-noise ratio (SNR).

    The general flow involves:
    1.  **Data Extraction**: Relevant simulation parameters such as satellite positions,
        detector characteristics (pointing, FOV, integration time, aperture area),
        target properties (positions, sizes, albedo), and celestial body positions
        (e.g., Sun) are extracted from the `sim_data` dictionary. If a `mask` is 
        provided, only the specified subset of detectors is extracted.
    2.  **Background Flux Calculation**: Utilizes `radiometry_calcs.fluxes` to determine
        the background contributions from the Sun, space, and sky based on the
        detector's filter.
    3.  **Visibility Determination**: For the active satellites, the function calculates
        the angular separation between each detector's pointing vector and the
        vectors to all fixed targets using vectorized broadcasting. A 2D mask is then
        applied to identify only those target-detector pairs that fall within the
        respective fields of view (FOV).
    4.  **Detector Flux Calculation**: For the visible targets, the incident flux
        on the detector is calculated using the `lambertian.lambertiansphere` model
        in a vectorized manner. This calculation considers the Sun's position, the 
        target's position, albedo, radius, and the solar flux.
    5.  **SNR Calculation**: Based on the calculated detector flux, the integration time,
        and the detector's aperture area, the signal and noise levels are determined
        simultaneously for all detections. The signal-to-noise ratio (SNR) is then 
        computed from these values.
    6.  **Output/Storage**: The calculated signal, noise, and SNR for each detected
        target are returned in a dictionary. If `print_output` is enabled, they are
        also printed to the console.

    Args:
        sim_data (dict): The main simulation data dictionary.
        print_output (int): If > 0, prints detailed results for each detection.
        mask (np.ndarray, optional): A boolean mask of length equal to the total number 
                                     of satellites, indicating which detectors to process.
                                     If None, all detectors are processed.

    Returns:
        dict: A dictionary containing simulation 'time' and arrays for 'sat_indices', 
              'target_indices', 'signal', 'noise', and 'snr'.

    TODO: we are using the same filter for all the detectors!
    """
    # 1. Data Extraction and Masking
    num_sats = sim_data.counts.get('satellites', 0)
    num_obs = sim_data.counts.get('observatories', 0)
    num_total_detectors = len(sim_data.detector.filt)
    sim_time = sim_data.time

    if mask is None:
        mask = np.ones(num_total_detectors, dtype=bool)
    
    active_indices = np.where(mask)[0]

    # Initialize empty result structure
    results = {
        'time': sim_time,
        'sat_indices': np.array([], dtype=int),
        'target_indices': np.array([], dtype=int),
        'signal': np.array([], dtype=float),
        'noise': np.array([], dtype=float),
        'snr': np.array([], dtype=float)
    }

    if len(active_indices) == 0:
        return results

    # Subset observer positions and detector properties cleanly
    satpositions = sim_data.get_detector_positions(mask)
    det_subset = sim_data.detector.subset(mask)

    detectorVect = det_subset.pointing
    fovs = det_subset.fov
    integrationTime = det_subset.integrationTime
    apertureArea = det_subset.apertureArea
    pixelOmega = det_subset.pixelOmega
    qe = det_subset.qe
    
    # Target and celestial data
    targets = sim_data.fixedpoints.position
    sunVect = sim_data.celestial.position[0]
    albedo = sim_data.fixedpoints.albedo
    radius = sim_data.fixedpoints.size / 2

    # 2. Background Flux Calculation
    # Assuming same filter for active group (as per existing TODO)
    filter_name = sim_data.detector.filt[active_indices[0]]
    sun, space, sky = radiometry_calcs.fluxes(filter_name)

    if print_output:
        print(f'Using filter: {filter_name}')
        print(f'sun, space, sky \n    {sun:.3e}, {space:.3e}, {sky:.3e}')

    # 3. Visibility Determination (Vectorized Geometry)
    # toTargets shape: (num_active_sats, num_targets, 3)
    toTargets = targets[np.newaxis, :, :] - satpositions[:, np.newaxis, :]
    
    # Calculate norms for normalization and distance
    norms_toTargets = np.linalg.norm(toTargets, axis=2)
    norms_toTargets_safe = np.where(norms_toTargets == 0, 1.0, norms_toTargets)
    
    # Normalized vectors from sats to targets
    normalized_toTargets = toTargets / norms_toTargets_safe[:, :, np.newaxis]
    
    # Vectorized dot product for angles: (num_active_sats, num_targets)
    dot_products = np.einsum('si,sti->st', detectorVect, normalized_toTargets)
        # This line uses NumPy's Einstein Summation (einsum) to perform a
        # high-performance, vectorized dot   product between the satellite
        # pointing vectors and the vectors to every target.
    angles = np.arccos(np.clip(dot_products, -1.0, 1.0))

    # Apply FOV mask: (num_active_sats, num_targets)
    # fov is the full field-of-view diameter; a target is visible
    # when its angular offset from boresight is less than fov / 2.
    visible_mask = angles < (fovs[:, np.newaxis] / 2)
    sat_hit_idx, target_hit_idx = np.where(visible_mask)
    
    if len(sat_hit_idx) == 0:
        if print_output:
            print("No targets visible to active detectors.")
        return results

    # 3.5 Occlusion Filter (Earth Limb & Horizon)
    hit_obs_pos = satpositions[sat_hit_idx]
    hit_toTargets_pre = toTargets[sat_hit_idx, target_hit_idx]
    hit_norms_pre = norms_toTargets[sat_hit_idx, target_hit_idx]
    
    u_hit_toTargets = hit_toTargets_pre / hit_norms_pre[:, np.newaxis]
    obs_norms = np.linalg.norm(hit_obs_pos, axis=1)
    valid_obs = obs_norms > 1e-9
    
    u_nadir = np.zeros_like(hit_obs_pos)
    u_nadir[valid_obs] = -hit_obs_pos[valid_obs] / obs_norms[valid_obs, np.newaxis]
    
    cos_angle_to_nadir = np.sum(u_hit_toTargets * u_nadir, axis=1)
    angle_to_nadir = np.arccos(np.clip(cos_angle_to_nadir, -1.0, 1.0))
    
    hit_earthEx = det_subset.earthEx[sat_hit_idx]
    hit_categories = np.array(det_subset.category)[sat_hit_idx]
    
    obs_is_ground = (hit_categories == 'observatories')

    # Ground observatory horizon check: angle from zenith must be <= (90deg - min_elevation)
    angle_to_zenith = np.pi - angle_to_nadir
    max_zenith = np.pi / 2.0 - hit_earthEx
    horizon_occluded = obs_is_ground & valid_obs & (angle_to_zenith > max_zenith)

    # Spacecraft Earth limb check: angle to Earth center must be >= (apparent_radius + earthEx)
    safe_obs_norms = np.where(obs_norms > EARTH_RADIUS, obs_norms, EARTH_RADIUS)
    apparent_radius = np.where(
        obs_norms > EARTH_RADIUS,
        np.arcsin(EARTH_RADIUS / safe_obs_norms),
        np.pi / 2.0
    )
    limb_occluded = (~obs_is_ground) & valid_obs & (angle_to_nadir < (apparent_radius + hit_earthEx))

    is_occluded = horizon_occluded | limb_occluded

    # Filter out occluded hits
    clear_hits = ~is_occluded
    sat_hit_idx = sat_hit_idx[clear_hits]
    target_hit_idx = target_hit_idx[clear_hits]
    
    if len(sat_hit_idx) == 0:
        if print_output:
            print("No targets visible (all occluded) to active detectors.")
        return results

    # 4. Detector Flux Calculation (Vectorized Radiometry)
    # Subset data for specific detector-target pairs (hits)
    hit_toTargets = toTargets[sat_hit_idx, target_hit_idx]
    hit_norms = norms_toTargets[sat_hit_idx, target_hit_idx]
    hit_albedo = albedo[target_hit_idx]
    hit_radius = radius[target_hit_idx]
    
    # Lambertian Phase Angle Calculation
    vec_from_sphere_to_observer = -hit_toTargets
    vec_from_sphere_to_light = np.tile(sunVect, (len(sat_hit_idx), 1))
    
    angle_light_observer = includedAngle(vec_from_sphere_to_light, vec_from_sphere_to_observer)

    target_brightness = lambertiansphere(
        angle_light_observer,
        hit_albedo,
        hit_radius,
        sun,
        debug=0 
    )

    detectorFlux = target_brightness / (np.pi * hit_norms ** 2)

    # 5. SNR Calculation
    hit_itime = integrationTime[sat_hit_idx]
    hit_aper = apertureArea[sat_hit_idx]
    hit_qe = qe[sat_hit_idx]
    hit_omega = pixelOmega[sat_hit_idx]
    hit_photoEff = sim_data.detector.photoEff[mask][sat_hit_idx]

    signal = detectorFlux * hit_itime * hit_aper * hit_qe * hit_photoEff
    noise = np.sqrt(space * hit_itime * hit_aper * hit_omega * hit_qe)
    snr = signal / noise

    # 6. Output/Storage
    if print_output:
        for k in range(len(sat_hit_idx)):
            i = active_indices[sat_hit_idx[k]]
            j = target_hit_idx[k]
            print(f"Sat {i} (itime={hit_itime[k]:.1f}s) detects Target {j}:")
            print(f"  Signal: {signal[k]:12.3e} | Noise: {noise[k]:12.3e} | SNR: {snr[k]:12.3e}")

    results.update({
        'sat_indices': active_indices[sat_hit_idx],
        'target_indices': target_hit_idx,
        'signal': signal,
        'noise': noise,
        'snr': snr
    })

    return results



\end{lstlisting}



\chapter{radiometry\_data.py: Filter Spectral Data}
\label{sec:radiometrydata}
\index{radiometry\_data.py}

\section*{Overview}
Defines FILTER\_DATA dictionary containing astronomical filter zero-points, sky/space background, and solar magnitudes.

\section*{Literate Source Code}
The source code of \texttt{radiometry\_data.py} is documented below. Section label: \ref{sec:radiometrydata}.

\begin{lstlisting}[language=Python,caption={radiometry\_data.py Source Code},label={code:radiometrydata}]
import scipy.constants as const

# --- Physical Constants ---
AU_M = const.au  # Astronomical Unit in meters
RSUN_M = 6.957e8 # Radius of the Sun in meters

# --- Radiometric Data ---

# Data is compiled for standard astronomical filters.
#
# Magnitudes are apparent magnitudes in the AB or Vega system.
# sun is the apparent magnitude of the sun in the given filter band.
# Sky brightness is for a dark site at new moon, in magnitudes per square arcsecond.
#        For space-based IR  telescopes, this is the zodiacal + telescope background.
#        For ground-based IR, it's dominated by thermal emission.
# Wavelengths and bandwidths are in nanometers (nm).
# Zero point is the photon flux for a 0-magnitude object in photons per second per square meter.



FILTER_DATA = {
    # Johnson-Cousins UBVRI Filters
    
    'U': {
        'sun': -26.03,
        'sky': 22.0,
        'space': 22.9,
        'central_wavelength': 365.0,
        'bandwidth': 66.0,
        'zero_point': 4.96e9,
    },
    'B': {
        'sun': -26.13,
        'sky': 22.7,
        'space': 23.3,
        'central_wavelength': 445.0,
        'bandwidth': 94.0,
        'zero_point': 1.36e10,
    },
    'V': {
        'sun': -26.78,
        'sky': 21.9,  # Mathematia was assuming 21.61, pretty close actually!
        'space': 23,
        'central_wavelength': 551.0,
        'bandwidth': 88.0,
        'zero_point': 8.79e9,
    },
    'R': {
        'sun': -27.11,
        'sky': 21.0,
        'space': 22.0,
        'central_wavelength': 658.0,
        'bandwidth': 138.0,
        'zero_point': 9.74e9,
    },
    'I': {
        'sun': -27.47,
        'sky': 20.0,
        'space': 21.2,
        'central_wavelength': 806.0,
        'bandwidth': 149.0,
        'zero_point': 7.11e9,
    },
    # Near-Infrared JHK Filters
    'J': {
        'sun': -27.91,
        'sky': 16.0,
        'space': 21.5,
        'central_wavelength': 1235.0,
        'bandwidth': 162.0,
        'zero_point': 3.15e9,
    },
    'H': {
        'sun': -28.25,
        'sky': 14.0,
        'space': 20.5,
        'central_wavelength': 1646.0,
        'bandwidth': 251.0,
        'zero_point': 2.35e9,
    },
    'K': {
        'sun': -28.29,
        'sky': 13.0,
        'space': 19.5,
        'central_wavelength': 2159.0,
        'bandwidth': 262.0,
        'zero_point': 1.17e9,
    },
    # SDSS Filters
    'g': {
        'sun': -26.47,
        'sky': 21.8,
        'space': 23.0,
        'central_wavelength': 477.0,
        'bandwidth': 137.9,
        'zero_point': 1.582e10,
    },
    'r': {
        'sun': -26.93,
        'sky': 20.8,
        'space': 22.2,
        'central_wavelength': 623.1,
        'bandwidth': 138.2,
        'zero_point': 1.214e10,
    },
    'i': {
        'sun': -27.05,
        'sky': 20.2,
        'space': 21.5,
        'central_wavelength': 762.5,
        'bandwidth': 153.5,
        'zero_point': 1.103e10,
    },
    'z': {
        'sun': -27.07,
        'sky': 19.0,
        'space': 20.8,
        'central_wavelength': 913.4,
        'bandwidth': 140.9,
        'zero_point': 8.455e9,
    },
    # Ground-Based Mid-IR Filters
    'L': {
        'sun': None,
        'sky': 3.5,
        'central_wavelength': 3500.0,
        'bandwidth': 600.0,
        'zero_point': 9.4e9,
    },
    'M': {
        'sun': None,
        'sky': 0.0,
        'central_wavelength': 4800.0,
        'bandwidth': 600.0,
        'zero_point': 6.8e9,
    },
    'N': {
        'sun': None,
        'sky': -6.0,
        'central_wavelength': 10200.0,
        'bandwidth': 5000.0,
        'zero_point': 2.7e10,
    },
    # JWST MIRI Filters
    'F560W': { 'sun': None,
               'sky': 24.4,
               'central_wavelength': 5600.0,
               'bandwidth': 1000.0,
               'zero_point': 9.78e9 },
    'F770W': { 'sun': None,
               'sky': 24.4,
               'central_wavelength': 7700.0,
               'bandwidth': 1950.0, 
               'zero_point': 1.39e10 },
    'F1000W': { 'sun': None,
                'sky': 23.6,
                'central_wavelength': 10000.0,
                'bandwidth': 1800.0, 
                'zero_point': 9.86e9 },
    'F1500W': { 'sun': None,
                'sky': 21.9,
                'central_wavelength': 15000.0,
                'bandwidth': 2940.0, 
                'zero_point': 1.08e10 },
    'F2550W': { 'sun': None,
                'sky': 19.4,
                'central_wavelength': 25500.0,
                'bandwidth': 4050.0,
                'zero_point': 8.70e9 }
}

\end{lstlisting}



\chapter{radiometry\_calcs.py: Radiometric Calculations \& Blackbody}
\label{sec:radiometrycalcs}
\index{radiometry\_calcs.py}

\section*{Overview}
Computes photon fluxes, magnitude conversions, and Planck blackbody spectral radiance.

\section*{Literate Source Code}
The source code of \texttt{radiometry\_calcs.py} is documented below. Section label: \ref{sec:radiometrycalcs}.

\begin{lstlisting}[language=Python,caption={radiometry\_calcs.py Source Code},label={code:radiometrycalcs}]
import math
import scipy.integrate
import scipy.constants as const
import numpy as np
import plotly.graph_objects as go
from radiometry_data import *
from constants import ARCSEC

def fluxes(band):
    """
    uses the FILTER_DATA table from radiometry_data.py for data
    Looks up in formation based on the argument band, which is 
    usually something like an astronomical band... U, B, V, etc.
    It returns three numbers:
    sun which is the solar flux at earth in photons/s/m^2
    sky which is the sky brightness at earth in p/s/steradian/m^2
    space, sky brightness in space in p/s/steradian/m^2
    """
    x = FILTER_DATA[band]
    zp = x['zero_point']
    sun_mag = x.get('sun')
    sun = amag(sun_mag) * zp if sun_mag is not None else 0.0

    space_mag = x.get('space', x.get('sky'))
    space = amag(space_mag) * zp / (ARCSEC**2) if space_mag is not None else 0.0

    sky_mag = x.get('sky')
    sky = amag(sky_mag) * zp / (ARCSEC**2) if sky_mag is not None else 0.0

    return (sun, space, sky)
    

def mag(x: float) -> float:
    """
    Calculates a magnitude value from a linear ratio.
    Uses the formula: magnitude = -2.5 * log10(ratio)
    """
    if x <= 0:
        raise ValueError("Input for mag() must be positive.")
    return -2.5 * math.log10(x)

def amag(x: float) -> float:
    """
    Calculates the linear ratio from a magnitude value.
    This is the inverse of the mag() function.
    Uses the formula: ratio = 10**(-0.4 * magnitude)
    """
    return 10**(-0.4 * x)


def _planck_law(wav_m: float, temp_k: float) -> float:
    """
    Helper function for Planck's law for spectral radiance.
    Args:
        wav_m: Wavelength in meters.
        temp_k: Temperature in Kelvin.
    Returns:
        Spectral radiance in W / (m^2 * sr * m).
    """
    if wav_m <= 0:
        return 0.0
    exponent = (const.h * const.c) / (wav_m * const.k * temp_k)
    if exponent > 700:
        return 0.0
    numerator = 2.0 * const.h * const.c**2
    denominator = (wav_m**5) * (math.expm1(exponent))
    if denominator == 0:
        return 0.0
    return numerator / denominator

def blackbody_flux(
    temperature: float,
    lambda_short: float,
    lambda_long: float
) -> float:
    """
    Numerically computes the integrated spectral radiance of
    a blackbody over a given wavelength band.

    Requires the scipy library.

    Args:
        temperature: The temperature of the blackbody in Kelvin.
        lambda_short: The short wavelength of the band in microns.
        lambda_long: The long wavelength of the band in microns.

    Returns:
        The integrated spectral radiance in units of
        Watts / (m^2 * steradian).
    """
    if lambda_short >= lambda_long:
        raise ValueError("lambda_short must be less than lambda_long.")
    lambda_short_m = lambda_short * 1e-6
    lambda_long_m = lambda_long * 1e-6
    integrated_radiance, _ = scipy.integrate.quad(
        lambda wav: _planck_law(wav, temperature),
        lambda_short_m,
        lambda_long_m
    )
    return integrated_radiance

def stefan_boltzmann_law(temperature: float) -> float:
    """
    Calculates the total power radiated per unit area by a
    blackbody using the Stefan-Boltzmann law.

    Args:
        temperature: The temperature of the blackbody in Kelvin.

    Returns:
        The total radiated power per unit area in W / m^2.
    """
    return const.sigma * (temperature**4)

def plot_blackbody_spectrum(temperature: float):
    """
    Plots the spectral radiance of a blackbody from
    0.5 to 30 microns.

    Args:
        temperature: The temperature of the blackbody in Kelvin.
    """
    wavelengths_microns = np.geomspace(0.5, 30, 500)
    wavelengths_m = wavelengths_microns * 1e-6
    spectral_radiance = [_planck_law(w, temperature) * 1e-6 for w in wavelengths_m]
    fig = go.Figure()
    fig.add_trace(go.Scatter(
        x=wavelengths_microns,
        y=spectral_radiance,
        mode='lines',
        name=f'{temperature} K Blackbody'
    ))
    fig.update_layout(
        title_text=f'Blackbody Spectrum ({temperature} K)',
        xaxis_title="Wavelength (microns)",
        yaxis_title="Spectral Radiance (W / m^2 / sr / micron)",
        xaxis_type="log",
        yaxis_type="log",
        template="plotly_white"
    )
    fig.show()

def plot_blackbody_spectrum_visible_nir(temperature: float):
    """
    Plots the spectral radiance of a blackbody from
    0.1 to 1 micron.

    Args:
        temperature: The temperature of the blackbody in Kelvin.
    """
    wavelengths_microns = np.linspace(0.1, 1.0, 500)
    wavelengths_m = wavelengths_microns * 1e-6
    spectral_radiance = [_planck_law(w, temperature) * 1e-6 for w in wavelengths_m]
    fig = go.Figure()
    fig.add_trace(go.Scatter(
        x=wavelengths_microns,
        y=spectral_radiance,
        mode='lines',
        name=f'{temperature} K Blackbody'
    ))
    fig.update_layout(
        title_text=f'Blackbody Spectrum ({temperature} K) - Visible/NIR',
        xaxis_title="Wavelength (microns)",
        yaxis_title="Spectral Radiance (W / m^2 / sr / micron)",
        template="plotly_white"
    )
    fig.show()



\end{lstlisting}



\chapter{lambertian.py: Lambertian Sphere Brightness}
\label{sec:lambertian}
\index{lambertian.py}

\section*{Overview}
Calculates phase-angle dependent reflected flux from Lambertian target spheres.

\section*{Literate Source Code}
The source code of \texttt{lambertian.py} is documented below. Section label: \ref{sec:lambertian}.

\begin{lstlisting}[language=Python,caption={lambertian.py Source Code},label={code:lambertian}]
import numpy as np

def lambertiansphere(
    angle_light_observer: np.ndarray,
    albedo: np.ndarray,
    radius: np.ndarray,
    base_brightness,
    debug: int = 0
) -> np.ndarray:
    """
    Calculates the emitted brightness from multiple lambertian spheres
    based on phase angle. This function computes the brightness emitted
    *from the sphere's surface* in the direction of the observer.
    The final apparent brightness at the observer's location must be
    calculated by dividing this result by pi * (distance_to_observer)^2
    to satisfy conservation of energy for a diffuse sphere.

    Args:
        angle_light_observer: A 1D NumPy array of shape (N,) with the
            phase angle in radians for each sphere. This is the angle
            between the light source and the observer as seen from
            the sphere's center.
        albedo: A 1D NumPy array of shape (N,) with the fraction of
            incident light that is reflected for each sphere (0.0 to 1.0).
        radius: A 1D NumPy array of shape (N,) with the radius of each
            sphere in meters.
        base_brightness: A 1D NumPy array of shape (N,) with the incident
            flux or brightness of the light source at each sphere's
            location.
        debug: An optional integer. If set to 1, a table of inputs and
            outputs will be printed before the function returns.
            Defaults to 0.

    Returns:
        A 1D NumPy array of shape (N,) containing the brightness emitted
        from each sphere's surface (e.g., in Watts per steradian per
        square meter). To get apparent brightness at the observer,
        divide by pi * (distance_to_observer)^2.
    """
    if not np.all((albedo >= 0.0) & (albedo <= 1.0)):
        raise ValueError("All albedo values must be between 0.0 and 1.0.")
    if np.any(radius < 0):
        raise ValueError("Radius cannot be negative.")
    if (not isinstance(angle_light_observer, np.ndarray) or
            angle_light_observer.ndim != 1):
        raise ValueError("angle_light_observer must be a 1D NumPy array.")

    # The phase angle is physically constrained to be between 0 and pi.
    # We clip the value to handle out-of-range inputs gracefully.
    alpha = np.clip(angle_light_observer, 0, np.pi)

    term1 = np.sin(alpha)
    term2 = (np.pi - alpha) * np.cos(alpha)
    phase_function_value = (2 / (3 * np.pi)) * (term1 + term2)

    cross_sectional_area = np.pi * (radius ** 2)

    effective_cross_section = (
        albedo *
        cross_sectional_area *
        phase_function_value
    )

    # The apparent brightness is the incident brightness multiplied
    # by the effective reflecting area.
    # The division by distance squared will be handled by the caller,
    # as instructed by the user, to accommodate the 4*pi*r^2 factor.
    apparent_brightness = (base_brightness * effective_cross_section)

    if debug == 1:
        print("\n--- Debug Info: lambertiansphere ---")
        print("Effective Cross Section: ", effective_cross_section)
        num_spheres = len(albedo)
        h1 = f"{'Index':<5} {'Phase Angle (rad)':<18} {'Albedo':<10}"
        h2 = f"{'Radius (m)':<12} {'Base Brightness':<18} {'Emitted Brightness':<22}"
        header = h1 + " " + h2
        print(header)
        print("-" * len(header))

        display_limit = 5
        for i in range(num_spheres):
            if (num_spheres > 2 * display_limit and
                    display_limit <= i < num_spheres - display_limit):
                if i == display_limit:
                    p1 = f"{'.':<5} {'...':<18} {'...':<10}"
                    p2 = f"{'...':<12} {'...':<18} {'...':<22}"
                    print(p1 + " " + p2)
                continue

            brightness_val = (
                base_brightness[i]
                if np.ndim(base_brightness) > 0
                else base_brightness
            )
            r1 = f"{i:<5} {angle_light_observer[i]:<18.4e} {albedo[i]:<10.4f}"
            r2 = f"{radius[i]:<12.4e} {brightness_val:<18.4e}"
            r3 = f"{apparent_brightness[i]:<22.4e}"
            print(f"{r1} {r2} {r3}")
        print("-----------------------------------\n")
    elif debug == 2:
        print("\n--- Detailed Debug Info: lambertiansphere ---")
        num_spheres = len(albedo)
        header = (
            f"{'Index':<5} {'Input Angle':<15} {'Clipped Alpha':<15} "
            f"{'Phase Func Val':<15} {'Cross Sect Area':<15} "
            f"{'Effective CS':<15} {'Albedo':<10} {'Radius (m)':<12} "
            f"{'Base Brightness':<18} {'Emitted Brightness':<22}"
        )
        print(header)
        print("-" * len(header))

        display_limit = 5
        for i in range(num_spheres):
            if (num_spheres > 2 * display_limit and
                    display_limit <= i < num_spheres - display_limit):
                if i == display_limit:
                    dots = ["..."] * 9
                    print(f"{'.':<5} " + " ".join([f"{d:<15}" for d in dots]))
                continue

            brightness_val = (
                base_brightness[i]
                if np.ndim(base_brightness) > 0
                else base_brightness
            )
            print(
                f"{i:<5} {angle_light_observer[i]:<15.4e} {alpha[i]:<15.4e} "
                f"{phase_function_value[i]:<15.4e} "
                f"{cross_sectional_area[i]:<15.4e} "
                f"{effective_cross_section[i]:<15.4e} {albedo[i]:<10.4f} "
                f"{radius[i]:<12.4e} {brightness_val:<18.4e} "
                f"{apparent_brightness[i]:<22.4e}"
            )
        print("-------------------------------------------\n")

    return apparent_brightness


def simple_lambertian(
    diameter: float,
    distance: float,
    albedo: float,
    angle: float,
    base_brightness: float
) -> float:
    """
    Calculates the apparent brightness of a Lambertian sphere.
    See Also lambertiansphere, which works with complete vectors
    of objects in a more code friendly way!

    This function computes the apparent brightness of a diffusely
    reflecting sphere based on its physical properties, viewing
    geometry, and a given base incident brightness. It simplifies
    the calculation by taking the phase angle directly, rather than
    calculating it from vectors.

    Args:
        diameter: The diameter of the sphere in meters.
        distance: The distance from the sphere to the observer
            in meters.
        albedo: The fraction of incident light that is
            reflected (0.0 to 1.0).
        angle: The phase angle in radians. This is the angle
            between the light source and the observer as seen
            from the sphere's center (expected to be between 0 and pi).
        base_brightness: The incident flux or brightness of the
            light source at the sphere's location (e.g., in
            Watts per square meter or photons / s / m^2).

    Returns:
        The apparent brightness of the sphere as observed from
        the specified distance (e.g., in Watts per square meter
        or photons / s / m^2).
    """
    if not 0.0 <= albedo <= 1.0:
        raise ValueError("Albedo must be between 0.0 and 1.0.")
    if diameter < 0:
        raise ValueError("Diameter cannot be negative.")
    if distance <= 0:
        raise ValueError("Distance must be positive.")

    # Calculate radius from diameter
    radius = diameter / 2.0

    # The phase angle is physically constrained to be between 0 and pi.
    # We clip the value to handle out-of-range inputs gracefully.
    angle = np.clip(angle, 0, np.pi)

    # This is the phase function for a Lambertian sphere.
    # It describes how the brightness changes with the phase angle.
    term1 = np.sin(angle)
    term2 = (np.pi - angle) * np.cos(angle)
    phase_function_value = (2 / (3 * np.pi)) * (term1 + term2)

    # Cross-sectional area of the sphere.
    cross_sectional_area = np.pi * (radius ** 2)

    # The effective cross-section combines the sphere's size,
    # reflectivity (albedo), and the phase function.
    effective_cross_section = (
        albedo *
        cross_sectional_area *
        phase_function_value
    )

    # The apparent brightness is the incident brightness multiplied
    # by the effective reflecting area, with the resulting light
    # distributed according to the inverse square law.
    apparent_brightness = (
        (base_brightness * effective_cross_section) /
        (np.pi * distance ** 2)
    )

    return apparent_brightness


def includedAngle(
    vectors1: np.ndarray,
    vectors2: np.ndarray
) -> np.ndarray:
    """
    Calculates the included angle in radians between corresponding vectors
    in two input NumPy arrays.

    Args:
        vectors1 (np.ndarray): A NumPy array of shape (N, 3) representing the
                               first set of vectors.
        vectors2 (np.ndarray): A NumPy array of shape (N, 3) representing the
                               second set of vectors.

    Returns:
        np.ndarray: A 1D NumPy array of shape (N,) containing the included
                    angle in radians for each corresponding pair of vectors.
    """
    if vectors1.shape != vectors2.shape:
        raise ValueError("Input arrays must have the same shape.")
    if vectors1.shape[1] != 3:
        raise ValueError("Input vectors must be 3-dimensional.")

    from minimalsimulation import safe_normalize
    unit_vectors1 = safe_normalize(vectors1, axis=1)
    unit_vectors2 = safe_normalize(vectors2, axis=1)

    # Calculate the dot product of the normalized vectors
    dot_product = np.einsum('ij,ij->i', unit_vectors1, unit_vectors2)

    # Clip values to ensure they are within the valid range for arccos (-1 to 1)
    dot_product = np.clip(dot_product, -1.0, 1.0)

    # Calculate the angle in radians
    angles = np.arccos(dot_product)

    return angles

\end{lstlisting}



\chapter{dataHandling.py: Simulation Data Logging \& Analysis}
\label{sec:datahandling}
\index{dataHandling.py}

\section*{Overview}
Collects observation results into DataHandler, supporting Pandas DataFrame conversion, CSV/Parquet export, and gap-time analysis.

\section*{Literate Source Code}
The source code of \texttt{dataHandling.py} is documented below. Section label: \ref{sec:datahandling}.

\begin{lstlisting}[language=Python,caption={dataHandling.py Source Code},label={code:datahandling}]
import pandas as pd
import numpy as np
import plotly.graph_objects as go

class DataHandler:
    """
    A class to collect, manage, and export simulation results from VibeVolts.
    
    This handler facilitates collecting results from multiple calls to 
    `nextIntegration` (via `scandetectors`) and converting them into a 
    structured Pandas DataFrame for analysis, plotting, and export.

    Functions:
        - __init__(): Initializes an empty results collection.
        - add_results(results): Appends a dictionary of scan results to the collection.
        - get_dataframe(): Consolidates all results into a single Pandas DataFrame.
        - save_to_csv(filename): Exports the collected data to a CSV file.
        - save_to_parquet(filename): Exports the collected data to a Parquet file.
            import pandas as pd
            df = pd.read_parquet('my_data.parquet')
        - clear(): Resets the handler by removing all collected data.
    """

    def __init__(self):
        """Initializes an empty list to store result dictionaries."""
        self.results_list = []
        self._master_df = None

    def add_results(self, results: dict):
        """
        Appends a result dictionary from nextIntegration to the collection.
        
        Args:
            results (dict): The dictionary returned by nextIntegration/scandetectors.
                            Expected to contain 'time', 'sat_indices', 
                            'target_indices', 'signal', 'noise', and 'snr'.
        """
        # We only add if there are actual detections to keep the DataFrame concise.
        # If you need to track empty time steps, remove this check.
        if len(results['sat_indices']) > 0:
            # Converting to DataFrame immediately ensures consistent types
            df_step = pd.DataFrame(results)
            self.results_list.append(df_step)
            # Reset master_df so it's recalculated on next access
            self._master_df = None

    def get_dataframe(self) -> pd.DataFrame:
        """
        Combines all collected results into a single Pandas DataFrame.
        
        Returns:
            pd.DataFrame: A DataFrame with columns for time, satellite indices,
                          target indices, signal, noise, and SNR.
                          Returns an empty DataFrame if no results have been added.
        """
        if not self.results_list:
            return pd.DataFrame()
            
        if self._master_df is None:
            self._master_df = pd.concat(self.results_list, ignore_index=True)
            
        return self._master_df

    def save_to_csv(self, filename: str):
        """
        Saves the collected results to a CSV file.
        
        Args:
            filename (str): The path/name of the file to save (e.g., 'results.csv').
        """
        df = self.get_dataframe()
        if not df.empty:
            df.to_csv(filename, index=False)
            print(f"Results saved to {filename}")
        else:
            print("No data to save.")

    def save_to_parquet(self, filename: str):
        """
        Saves the collected results to a Parquet file (more efficient than CSV).
        
        Args:
            filename (str): The path/name of the file to save (e.g., 'results.parquet').
        """
        df = self.get_dataframe()
        if not df.empty:
            df.to_parquet(filename, index=False)
            print(f"Results saved to {filename}")
        else:
            print("No data to save.")

    def clear(self):
        """Clears all collected results to start a new run."""
        self.results_list = []
        self._master_df = None

    def calculate_gap_times(
        self, target_id: int = None, pooled: bool = False
    ) -> dict[int, np.ndarray] | np.ndarray:
        """
        Calculates the time gaps (in seconds) between consecutive observations of targets.

        If target_id is specified, returns a 1D numpy array of gap times (in seconds)
        for that target.
        If target_id is None and pooled is True, returns a single 1D numpy array containing
        the gap times of all targets concatenated.
        If target_id is None and pooled is False, returns a dictionary mapping target
        indices to their respective numpy arrays of gap times.

        Args:
            target_id (int, optional): The index of a specific target. Defaults to None.
            pooled (bool, optional): If True and target_id is None, returns a single
                                     concatenated array of all gap times. Defaults to False.

        Returns:
            np.ndarray or dict: Array of gap times or a dictionary mapping
                               target index -> numpy array of gap times.
        """
        df = self.get_dataframe()
        if df.empty:
            return np.array([], dtype=float) if target_id is not None else {}

        # Ensure sorted chronologically
        df_sorted = df.sort_values('time')

        if target_id is not None:
            target_df = df_sorted[df_sorted['target_indices'] == target_id]
            if len(target_df) < 2:
                return np.array([], dtype=float)
            time_diffs = target_df['time'].diff().dropna()
            return time_diffs.dt.total_seconds().values
        elif pooled:
            all_gaps = []
            for tid, group in df_sorted.groupby('target_indices'):
                if len(group) >= 2:
                    time_diffs = group['time'].diff().dropna()
                    all_gaps.extend(time_diffs.dt.total_seconds().tolist())
            return np.array(all_gaps, dtype=float)
        else:
            gaps_dict = {}
            for tid, group in df_sorted.groupby('target_indices'):
                if len(group) < 2:
                    gaps_dict[int(tid)] = np.array([], dtype=float)
                else:
                    time_diffs = group['time'].diff().dropna()
                    gaps_dict[int(tid)] = time_diffs.dt.total_seconds().values
            return gaps_dict

    def plot_gap_times_histogram(
        self,
        target_id: int = None,
        bins: int | str | None = "auto",
        show_plot: bool = True
    ) -> go.Figure:
        """
        Generates and optionally displays a Plotly histogram of the target gap times.

        If target_id is specified, plots the histogram of gap times for that target.
        If target_id is None, plots the histogram of pooled gap times across all targets.

        Args:
            target_id (int, optional): The index of a specific target. Defaults to None.
            bins (int or str, optional): Specification of histogram bins (e.g. 'auto', 10).
                                         Passed to Plotly's nbinsx. Defaults to 'auto'.
            show_plot (bool, optional): If True, displays the plot. Defaults to True.

        Returns:
            plotly.graph_objects.Figure: The generated Plotly figure object.
        """
        if target_id is not None:
            gaps = self.calculate_gap_times(target_id=target_id)
            title = f"Interobservation Gap Times for Target {target_id}"
        else:
            gaps = self.calculate_gap_times(pooled=True)
            title = "Pooled Interobservation Gap Times (All Targets)"

        if len(gaps) == 0:
            fig = go.Figure()
            fig.update_layout(
                title_text=f"{title} - No Gap Data Available",
                template="plotly_white"
            )
            if show_plot:
                fig.show()
            return fig

        fig = go.Figure()
        fig.add_trace(go.Histogram(
            x=gaps,
            nbinsx=bins if isinstance(bins, int) else None,
            marker=dict(
                color='#3B82F6',
                line=dict(color='white', width=1)
            ),
            opacity=0.85,
            name="Gap Times"
        ))

        fig.update_layout(
            title_text=title,
            xaxis_title="Gap Time (seconds)",
            yaxis_title="Observation Count",
            template="plotly_white",
            bargap=0.05,
            hovermode="x unified"
        )

        if show_plot:
            fig.show()

        return fig

\end{lstlisting}



\chapter{cadenceController.py: Discrete Event Scan Scheduler}
\label{sec:cadencecontroller}
\index{cadenceController.py}

\section*{Overview}
Groups detectors by integration time and schedules discrete integration events.

\section*{Literate Source Code}
The source code of \texttt{cadenceController.py} is documented below. Section label: \ref{sec:cadencecontroller}.

\begin{lstlisting}[language=Python,caption={cadenceController.py Source Code},label={code:cadencecontroller}]
import numpy as np
from datetime import timedelta
from propagation import propagate_satellites
from scandetectors import scandetectors
from celestialbodies import celestial_update
from minimalsimulation import CadenceGroup, CadenceState
from observatories import propagate_observatories


def initCadence(sim_data) -> None:
    """
    Initialises sim_data.cadenceStructure from detector integration times.

    Groups detectors that share the same integration time into CadenceGroup
    objects and stores them in a CadenceState on sim_data.cadenceStructure.
    All groups are initialised with scanNext equal to the current simulation
    time so that every group performs a scan on the first call to
    nextIntegration.

    Args:
        sim_data: The main SimulationState object.
    """
    integration_times = sim_data.detector.integrationTime
    unique_intervals = np.unique(integration_times)

    cadence_list = [
        CadenceGroup(
            scanInterval=float(interval),
            scanMask=(integration_times == interval),
            scanNext=sim_data.time,
        )
        for interval in unique_intervals
    ]

    sim_data.cadenceStructure = CadenceState(cadenceList=cadence_list)
    _update_next_schedule(sim_data)


def nextIntegration(sim_data, print_output: int = 0) -> dict:
    """
    Advances the simulation to the next scheduled scan and executes it.

    Steps:
    1. Advances sim_data.time to the earliest scheduled group event.
    2. Propagates all satellites and updates celestial body positions.
    3. Runs scandetectors for only the active group's masked detectors.
    4. Updates that group's scanNext by one scanInterval.
    5. Recomputes which group fires next.

    Args:
        sim_data:      The main SimulationState object.
        print_output:  If > 0, scandetectors prints per-detection details.

    Returns:
        dict: The detection results returned by scandetectors, containing
              'time', 'sat_indices', 'target_indices', 'signal', 'noise',
              and 'snr'.
    """
    cadence = sim_data.cadenceStructure

    # 1. Advance simulation time to the next scheduled event.
    sim_data.time = cadence.nextTime

    # 2. Propagate all satellites and celestial bodies to the new time.
    propagate_satellites(sim_data, sim_data.time)
    celestial_update(sim_data, sim_data.time)
    propagate_observatories(sim_data, sim_data.time)

    # 3. Run the scan for the active group's detector subset.
    group = cadence.cadenceList[cadence.nextGroup]
    results = scandetectors(sim_data, print_output=print_output,
                            mask=group.scanMask)

    # 4. Schedule this group's next scan.
    group.scanNext = sim_data.time + timedelta(seconds=group.scanInterval)

    # 5. Find the next overall event across all groups.
    _update_next_schedule(sim_data)

    return results


def _update_next_schedule(sim_data) -> None:
    """
    Scans all CadenceGroups and updates sim_data.cadenceStructure with
    the index and datetime of the earliest upcoming scan.

    Args:
        sim_data: The main SimulationState object.
    """
    cadence = sim_data.cadenceStructure
    next_times = [g.scanNext for g in cadence.cadenceList]
    min_idx = min(range(len(next_times)), key=lambda i: next_times[i])
    cadence.nextGroup = min_idx
    cadence.nextTime = cadence.cadenceList[min_idx].scanNext

\end{lstlisting}



\chapter{fibonacciSearch.py: Fibonacci Pointing Vector Resorting}
\label{sec:fibonaccisearch}
\index{fibonacciSearch.py}

\section*{Overview}
Generates Fibonacci spherical grids and re-orders vectors to minimize slew distances.

\section*{Literate Source Code}
The source code of \texttt{fibonacciSearch.py} is documented below. Section label: \ref{sec:fibonaccisearch}.

\begin{lstlisting}[language=Python,caption={fibonacciSearch.py Source Code},label={code:fibonaccisearch}]
# fibonacciSearch
# this contains the functions necessary to support a set of detectors
# that searches the sphere using a fibonacci grid

import numpy as np
import plotly.graph_objects as go

# def searchStruct(sim_data,detect):
#     '''
#     creates the data structure for each of the satellite detectors,
#     adds the structure to the detector
#     '''
#     theta = detect[:,IFOV_IDX]/2

#     # Calculate solid angle 
#     theta = fov / 2
#     solid_angle = 2 * np.pi * (1 - np.cos(theta))
    
#     # Calculate grid_points - blow things up by 0.25 for overlap
#     grid_points = int(4 * np.pi / solid_angle * 1.25)

#     # Generate and store the pointing sphere and place in ['pointing_sphers'][n]
#     generate_pointing_sphere(sim_data, grid_points)

def pointing_vectors(n: int) -> np.ndarray:
    """
    Generates n equally spaced points on a unit sphere using the Fibonacci lattice algorithm.

    Args:
        n: The number of points to generate.

    Returns:
        A NumPy array of shape (n, 3) for the Cartesian coordinates of the points.
    """
    if not isinstance(n, int) or n <= 0:
        raise ValueError("n must be a positive integer.")

    indices = np.arange(0, n, dtype=float) + 0.5
    z = 1 - 2 * indices / n
    radius_xy = np.sqrt(1 - z**2)
    golden_angle = np.pi * (3. - np.sqrt(5.))
    theta = golden_angle * indices
    x = radius_xy * np.cos(theta)
    y = radius_xy * np.sin(theta)

    unit_vectors = np.stack([x, y, z], axis=1)
    return unit_vectors

def resort_vectors_by_proximity(unit_vectors: np.ndarray) -> np.ndarray:
    """
    Resorts a list of vectors by making each subsequent vector the closest one
    in the remaining set to the previous one.

    Args:
        unit_vectors: A NumPy array of shape (n, 3) representing the vectors.

    Returns:
        A new NumPy array with the reordered vectors.
    """
    if unit_vectors.ndim != 2 or unit_vectors.shape[1] != 3:
        raise ValueError("unit_vectors array must have shape (n, 3).")

    n_vectors = unit_vectors.shape[0]
    if n_vectors == 0:
        return np.array([])

    remaining_indices = list(range(n_vectors))
    sorted_vectors = np.zeros_like(unit_vectors)
    
    # Start with the first vector
    current_index = remaining_indices.pop(0)
    sorted_vectors[0] = unit_vectors[current_index]
    
    for i in range(1, n_vectors):
        last_vector = sorted_vectors[i-1]
        
        # Vectorized distance computation to all remaining candidates
        rem_coords = unit_vectors[remaining_indices]
        dists = np.linalg.norm(rem_coords - last_vector, axis=1)
        
        # Find index of minimum distance
        best_idx_in_rem = np.argmin(dists)
        best_index = remaining_indices[best_idx_in_rem]

        sorted_vectors[i] = unit_vectors[best_index]
        remaining_indices.pop(best_idx_in_rem)

    return sorted_vectors


def plot_vectors_on_sphere(vectors: np.ndarray, title: str) -> go.Figure:
    """
    Creates a 3D plot of vectors on a sphere.

    Args:
        vectors: A NumPy array of shape (n, 3) representing the vectors.
        title: The title of the plot.

    Returns:
        A Plotly figure object.
    """
    if vectors.ndim != 2 or vectors.shape[1] != 3:
        raise ValueError("vectors array must have shape (n, 3).")

    fig = go.Figure()

    fig.add_trace(go.Scatter3d(
        x=vectors[:, 0], y=vectors[:, 1], z=vectors[:, 2],
        mode='markers',
        marker=dict(
            size=2,
            color='red',
            opacity=0.8
        ),
        name='Vectors'
    ))

    u_sphere = np.linspace(0, 2 * np.pi, 100)
    v_sphere = np.linspace(0, np.pi, 100)
    x_earth = 1.0 * np.outer(np.cos(u_sphere), np.sin(v_sphere))
    y_earth = 1.0 * np.outer(np.sin(u_sphere), np.sin(v_sphere))
    z_earth = 1.0 * np.outer(np.ones(np.size(u_sphere)), np.cos(v_sphere))
    fig.add_trace(go.Surface(
        x=x_earth, y=y_earth, z=z_earth,
        colorscale='Blues', showscale=False, opacity=0.5, name='Sphere'
    ))

    fig.update_layout(
        title=title,
        scene=dict(
            xaxis_title='X', yaxis_title='Y', zaxis_title='Z',
            aspectmode='data'
        ),
        margin=dict(r=20, b=10, l=10, t=40),
        legend_title_text='Objects'
    )

    return fig

def test_vector_resorting():
    """
    Tests the vector resorting and plots the Euclidean distance between subsequent vectors.
    """
    # 1. Generate 100 vectors
    vectors = pointing_vectors(100)

    # 2. Calculate distances without resorting
    diffs_unsorted = np.linalg.norm(np.diff(vectors, axis=0), axis=1)

    # 3. Resort the vectors
    sorted_vectors = resort_vectors_by_proximity(vectors)
    
    # 4. Calculate distances with resorting
    diffs_sorted = np.linalg.norm(np.diff(sorted_vectors, axis=0), axis=1)

    # 5. Plot the results
    fig = go.Figure()
    fig.add_trace(go.Scatter(y=diffs_unsorted, mode='lines', name='Unsorted'))
    fig.add_trace(go.Scatter(y=diffs_sorted, mode='lines', name='Resorted by Proximity'))
    fig.update_layout(
        title="Euclidean Distance Between Subsequent Vectors",
        xaxis_title="Vector Index",
        yaxis_title="Euclidean Distance"
    )
    fig.add_annotation(
        text="Generated by: test_vector_resorting",
        xref="paper", yref="paper",
        x=0.5, y=1.01,
        showarrow=False,
        font=dict(size=10, color="gray"),
        xanchor="center", yanchor="bottom"
    )
    return fig

if __name__ == '__main__':
    fig = test_vector_resorting()
    fig.show()

\end{lstlisting}



\chapter{generate\_log\_spherical\_points.py: Logarithmic Radial Distribution}
\label{sec:generatelogsphericalpoints}
\index{generate\_log\_spherical\_points.py}

\section*{Overview}
Generates 3D points with logarithmic radial and uniform angular distribution.

\section*{Literate Source Code}
The source code of \texttt{generate\_log\_spherical\_points.py} is documented below. Section label: \ref{sec:generatelogsphericalpoints}.

\begin{lstlisting}[language=Python,caption={generate\_log\_spherical\_points.py Source Code},label={code:generatelogsphericalpoints}]
import numpy as np
from datetime import datetime, timezone
from plotting_3d import plot_3d_scatter
from constants import GEO_RADIUS


def generate_log_spherical_points(
    num_points: int,
    inner_radius: float,
    outer_radius: float,
    seed: int = None
) -> np.ndarray:
    """
    Generates 3D points with logarithmic radial and uniform angular distribution.

    This function creates a point cloud where point distances from the origin are
    logarithmically spaced. On any given spherical shell, points are distributed
    uniformly using the Fibonacci lattice method.
    

    Args:
        num_points: The total number of points to generate.
        inner_radius: The minimum distance from the origin (must be positive).
        outer_radius: The maximum distance from the origin (must be >= inner_radius).
        seed: An optional integer to seed the random number generator for
              reproducible shuffling.

    Returns:
        A NumPy array of shape (num_points, 3) for the Cartesian
        coordinates.
    """
    # Input validation
    if not isinstance(num_points, int) or num_points <= 0:
        raise ValueError("num_points must be a positive integer.")
    if not (isinstance(inner_radius, (int, float)) and inner_radius > 0):
        raise ValueError("inner_radius must be a positive number.")
    if not (isinstance(outer_radius, (int, float)) and outer_radius >= inner_radius):
        raise ValueError("outer_radius must be a positive number and >= inner_radius.")

    # --- 1. Generate uniformly distributed unit vectors (Fibonacci Lattice) ---
    from fibonacciSearch import pointing_vectors
    unit_vectors = pointing_vectors(num_points)

    # --- 2. Generate logarithmically spaced radii ---
    # The start and stop parameters for logspace are exponents
    start_exp = np.log10(inner_radius)
    stop_exp = np.log10(outer_radius)
    radii = np.logspace(start_exp, stop_exp, num_points)

    # --- 3. Shuffle radii to decouple radius from spherical position ---
    # This prevents the smallest radii from always mapping to points near the
    # north pole, creating a more uniform visual distribution of densities.
    rng = np.random.default_rng(seed)
    rng.shuffle(radii)

    # --- 4. Scale unit vectors by radii using broadcasting ---
    # Reshape radii to (N, 1) to broadcast with (N, 3) unit_vectors
    points = unit_vectors * radii[:, np.newaxis]

    return points

if __name__ == '__main__':
    # --- Demo of Point Generation and Visualization ---
    # This provides a consistent demonstration entry point, matching vibevolts_demo.py
    NUM_POINTS = 100
    # Using the same radii as the main simulation for consistency
    INNER_RADIUS = 2000000
    OUTER_RADIUS = 2 * GEO_RADIUS
    OBJECT_SIZE = 1.0

    print(f"Generating {NUM_POINTS} points from radius {INNER_RADIUS} to {OUTER_RADIUS}...")
    generated_points = generate_log_spherical_points(
        num_points=NUM_POINTS,
        inner_radius=INNER_RADIUS,
        outer_radius=OUTER_RADIUS,
        seed=42
    )
    print("Point generation complete.\n")

    # Use the unified visualization function
    plot_3d_scatter(
        positions=generated_points,
        title="Fixed Points Distribution (from generate_log_spherical_points.py)",
        plot_time=datetime.now(timezone.utc),
        labels=[f"Point {i}" for i in range(len(generated_points))],
        marker_size=1,
        trace_name='Fixed Points'
    )

\end{lstlisting}



\chapter{plotting\_3d.py: Plotly 3D Position Renderer}
\label{sec:plotting3d}
\index{plotting\_3d.py}

\section*{Overview}
Renders 3D interactive Plotly figures for satellite, observatory, and target positions.

\section*{Literate Source Code}
The source code of \texttt{plotting\_3d.py} is documented below. Section label: \ref{sec:plotting3d}.

\begin{lstlisting}[language=Python,caption={plotting\_3d.py Source Code},label={code:plotting3d}]
import numpy as np
from datetime import datetime
from typing import List, Optional

import plotly.graph_objects as go
from astropy.time import Time
from astropy.coordinates import GCRS, EarthLocation
import astropy.units as u

from constants import EARTH_RADIUS

def plot_3d_scatter(
    positions: np.ndarray,
    title: str,
    plot_time: datetime,
    labels: Optional[List[str]] = None,
    marker_size: int = 1,
    trace_name: str = 'Points',
    marker_color: Optional[str] = None
) -> go.Figure:
    """
    Creates a 3D plot of object positions.
    """
    if positions.ndim != 2 or positions.shape[1] != 3:
        raise ValueError("positions array must have shape (n, 3).")

    fig = go.Figure()

    marker_dict = dict(
        size=marker_size,
        opacity=0.8
    )
    if marker_color:
        marker_dict['color'] = marker_color
    else:
        marker_dict['color'] = np.arange(len(positions))
        marker_dict['colorscale'] = 'Viridis'

    fig.add_trace(go.Scatter3d(
        x=positions[:, 0], y=positions[:, 1], z=positions[:, 2],
        mode='markers',
        marker=marker_dict,
        text=labels,
        hoverinfo='text' if labels else 'none',
        name=trace_name
    ))

    u_sphere = np.linspace(0, 2 * np.pi, 100)
    v_sphere = np.linspace(0, np.pi, 100)
    x_earth = EARTH_RADIUS * np.outer(np.cos(u_sphere), np.sin(v_sphere))
    y_earth = EARTH_RADIUS * np.outer(np.sin(u_sphere), np.sin(v_sphere))
    z_earth = EARTH_RADIUS * np.outer(np.ones(np.size(u_sphere)), np.cos(v_sphere))
    fig.add_trace(go.Surface(
        x=x_earth, y=y_earth, z=z_earth,
        colorscale='Blues', showscale=False, opacity=0.5, name='Earth'
    ))

    theta = np.linspace(0, 2 * np.pi, 100)
    x_eq = EARTH_RADIUS * np.cos(theta)
    y_eq = EARTH_RADIUS * np.sin(theta)
    z_eq = np.zeros_like(theta)
    fig.add_trace(go.Scatter3d(
        x=x_eq, y=y_eq, z=z_eq, mode='lines',
        line=dict(color='green', width=3), name='Equator'
    ))

    fig.add_trace(go.Scatter3d(
        x=[0], y=[0], z=[EARTH_RADIUS * 1.1], mode='text',
        text=['N'], textfont=dict(size=15, color='red'), name='North Pole'
    ))

    lat_es, lon_es = 33.92 * u.deg, -118.42 * u.deg
    el_segundo_loc = EarthLocation.from_geodetic(lon=lon_es, lat=lat_es)
    itrs_coords = el_segundo_loc.get_itrs(obstime=Time(plot_time))
    gcrs_coords = itrs_coords.transform_to(GCRS(obstime=Time(plot_time)))
    es_pos = gcrs_coords.cartesian.xyz.to(u.m).value * 1.05
    fig.add_trace(go.Scatter3d(
        x=[es_pos[0]], y=[es_pos[1]], z=[es_pos[2]], mode='text',
        text=['ES'], textfont=dict(size=15, color='yellow'), name='El Segundo'
    ))

    fig.update_layout(
        title=title,
        scene=dict(
            xaxis_title='X (m)', yaxis_title='Y (m)', zaxis_title='Z (m)',
            aspectmode='data'
        ),
        margin=dict(r=20, b=10, l=10, t=40),
        legend_title_text='Objects'
    )

    return fig

\end{lstlisting}



\chapter{plotting\_vectors.py: Plotly Pointing Vector Renderer}
\label{sec:plottingvectors}
\index{plotting\_vectors.py}

\section*{Overview}
Renders 3D interactive Plotly figures for sensor pointing vectors.

\section*{Literate Source Code}
The source code of \texttt{plotting\_vectors.py} is documented below. Section label: \ref{sec:plottingvectors}.

\begin{lstlisting}[language=Python,caption={plotting\_vectors.py Source Code},label={code:plottingvectors}]
import numpy as np
from datetime import datetime
from typing import Dict, Any

import plotly.graph_objects as go

from constants import EARTH_RADIUS

def plot_pointing_vectors(
    data_struct: Dict[str, Any],
    title: str,
    plot_time: datetime
) -> go.Figure:
    """
    Creates a 3D plot of satellites with pointing vectors.
    """
    sat_positions = data_struct.satellites.position.copy()
    sat_pointing = data_struct.detector.pointing.copy()
    num_sats = data_struct.counts.satellites

    if num_sats == 0:
        print("No satellites to plot.")
        return go.Figure()

    fig = go.Figure()

    fig.add_trace(go.Scatter3d(
        x=sat_positions[:, 0], y=sat_positions[:, 1], z=sat_positions[:, 2],
        mode='markers',
        marker=dict(size=5, color='blue', opacity=0.8),
        name='Satellites'
    ))

    vector_scale = 0.5 * EARTH_RADIUS
    for i in range(num_sats):
        start_point = sat_positions[i]
        pointing_vec = sat_pointing[i]
        norm = np.linalg.norm(pointing_vec)
        unit_vec = pointing_vec / norm if norm > 0 else np.array([0, 0, 0])
        end_point = start_point + unit_vec * vector_scale
        fig.add_trace(go.Scatter3d(
            x=[start_point[0], end_point[0]],
            y=[start_point[1], end_point[1]],
            z=[start_point[2], end_point[2]],
            mode='lines',
            line=dict(color='red', width=3),
            showlegend=False
        ))

    u_sphere = np.linspace(0, 2 * np.pi, 100)
    v_sphere = np.linspace(0, np.pi, 100)
    x_earth = EARTH_RADIUS * np.outer(np.cos(u_sphere), np.sin(v_sphere))
    y_earth = EARTH_RADIUS * np.outer(np.sin(u_sphere), np.sin(v_sphere))
    z_earth = EARTH_RADIUS * np.outer(np.ones(np.size(u_sphere)), np.cos(v_sphere))
    fig.add_trace(go.Surface(
        x=x_earth, y=y_earth, z=z_earth,
        colorscale='Blues', showscale=False, opacity=0.5, name='Earth'
    ))

    fig.update_layout(
        title=title,
        scene=dict(
            xaxis_title='X (m)', yaxis_title='Y (m)', zaxis_title='Z (m)',
            aspectmode='data'
        ),
        margin=dict(r=20, b=10, l=10, t=40)
    )

    return fig

\end{lstlisting}



\chapter{sim\_report.py: Simulation State Reporting}
\label{sec:simreport}
\index{sim\_report.py}

\section*{Overview}
Prints formatted state reports for satellites, observatories, and exclusion parameters.

\section*{Literate Source Code}
The source code of \texttt{sim\_report.py} is documented below. Section label: \ref{sec:simreport}.

\begin{lstlisting}[language=Python,caption={sim\_report.py Source Code},label={code:simreport}]
import numpy as np
from astropy.coordinates import EarthLocation, GCRS, AltAz, SkyCoord
from astropy.time import Time
import astropy.units as u

from constants import (
    ORBITAL_A_IDX,
    ORBITAL_E_IDX,
    ORBITAL_I_IDX,
    ORBITAL_RAAN_IDX,
    ORBITAL_ARGP_IDX,
    ORBITAL_M_IDX
)

def print_simulation_state(sim_data):
    """
    Prints the current state of the simulation in tabulated format.
    Includes satellites (orbital elements, pointing RA/Dec),
    observatories (lat/lon/alt, pointing Az/El), and exclusion parameters.
    """
    print("\n" + "="*80)
    print("SIMULATION STATE REPORT")
    print(f"Time: {sim_data.time}")
    print("="*80)

    # 1. Satellites
    print("\n--- SATELLITES ---")
    if 'satellites' in sim_data.counts and sim_data.counts.satellites > 0:
        num_sats = sim_data.counts.satellites
        print(f"{'Idx':<4} | {'Semi-Major (m)':<15} | {'Eccentricity':<12} | {'Inc (deg)':<9} | {'RAAN (deg)':<10} | {'Sensor Band':<11} | {'Pointing RA (deg)':<17} | {'Pointing Dec (deg)':<18}")
        print("-" * 115)
        
        # Build lookup array for detectors safely if detector is present
        has_detectors = (sim_data.detector is not None and len(sim_data.detector.filt) > 0)
        if has_detectors:
            det_cat = np.array(sim_data.detector.category)
            det_idx = sim_data.detector.asset_index
        else:
            det_cat, det_idx = np.array([]), np.array([])
        
        for i in range(num_sats):
            oe = sim_data.satellites.orbital_elements[i]
            a = oe[ORBITAL_A_IDX]
            e = oe[ORBITAL_E_IDX]
            inc = np.rad2deg(oe[ORBITAL_I_IDX])
            raan = np.rad2deg(oe[ORBITAL_RAAN_IDX])
            
            match = np.where((det_cat == 'satellites') & (det_idx == i))[0]
            if len(match) > 0:
                d_i = match[0]
                band = sim_data.detector.filt[d_i]
                pt = sim_data.detector.pointing[d_i]
                # Pointing vector is in GCRS
                ra = np.rad2deg(np.arctan2(pt[1], pt[0])) % 360
                dec = np.rad2deg(np.arcsin(pt[2]))
                band_str = band
                ra_str = f"{ra:.2f}"
                dec_str = f"{dec:.2f}"
            else:
                band_str = "None"
                ra_str = "N/A"
                dec_str = "N/A"
                
            print(f"{i:<4} | {a:<15.2f} | {e:<12.6f} | {inc:<9.2f} | {raan:<10.2f} | {band_str:<11} | {ra_str:<17} | {dec_str:<18}")
    else:
        print("No satellites in simulation.")

    # 2. Observatories
    print("\n--- OBSERVATORIES ---")
    if 'observatories' in sim_data.counts and sim_data.counts.observatories > 0:
        num_obs = sim_data.counts.observatories
        print(f"{'Idx':<4} | {'Lat (deg)':<9} | {'Lon (deg)':<10} | {'Alt (m)':<9} | {'Sensor Band':<11} | {'Pointing Az (deg)':<17} | {'Pointing El (deg)':<18}")
        print("-" * 95)
        
        if has_detectors:
            det_cat = np.array(sim_data.detector.category)
            det_idx = sim_data.detector.asset_index
        else:
            det_cat, det_idx = np.array([]), np.array([])
            
        obs_time = Time(sim_data.time)

        for i in range(num_obs):
            lat = sim_data.observatories.latitude[i]
            lon = sim_data.observatories.longitude[i]
            alt = sim_data.observatories.altitude[i]
            
            match = np.where((det_cat == 'observatories') & (det_idx == i))[0]
            if len(match) > 0:
                d_i = match[0]
                band = sim_data.detector.filt[d_i]
                pt = sim_data.detector.pointing[d_i]
                
                # Pointing vector is in GCRS. Convert to AltAz frame.
                loc = EarthLocation(lat=lat*u.deg, lon=lon*u.deg, height=alt*u.m)
                coord = SkyCoord(x=pt[0], y=pt[1], z=pt[2], representation_type='cartesian', frame=GCRS(obstime=obs_time))
                altaz = coord.transform_to(AltAz(obstime=obs_time, location=loc))
                
                az = altaz.az.deg
                el = altaz.alt.deg
                
                band_str = band
                az_str = f"{az:.2f}"
                el_str = f"{el:.2f}"
            else:
                band_str = "None"
                az_str = "N/A"
                el_str = "N/A"
                
            print(f"{i:<4} | {lat:<9.4f} | {lon:<10.4f} | {alt:<9.1f} | {band_str:<11} | {az_str:<17} | {el_str:<18}")
    else:
        print("No observatories in simulation.")

    # 3. Exclusion Parameters
    print("\n--- EXCLUSION PARAMETERS ---")
    if has_detectors:
        print(f"{'Category':<15} | {'Asset Idx':<9} | {'Solar Excl (deg)':<16} | {'Lunar Excl (deg)':<16} | {'Earth Excl (deg)':<16}")
        print("-" * 80)
        
        for d_i in range(len(sim_data.detector.filt)):
            cat = sim_data.detector.category[d_i]
            a_idx = sim_data.detector.asset_index[d_i]
            solar_ex = np.rad2deg(sim_data.detector.solarEx[d_i])
            lunar_ex = np.rad2deg(sim_data.detector.lunarEx[d_i])
            earth_ex = np.rad2deg(sim_data.detector.earthEx[d_i])
            
            print(f"{cat:<15} | {a_idx:<9} | {solar_ex:<16.2f} | {lunar_ex:<16.2f} | {earth_ex:<16.2f}")
    else:
        print("No detectors configured.")
        
    print("="*80 + "\n")

\end{lstlisting}



\chapter{sim\_check.py: Simulation Data Verification}
\label{sec:simcheck}
\index{sim\_check.py}

\section*{Overview}
Utility for checking keys, counts, and non-zero initialization of sim\_data.

\section*{Literate Source Code}
The source code of \texttt{sim\_check.py} is documented below. Section label: \ref{sec:simcheck}.

\begin{lstlisting}[language=Python,caption={sim\_check.py Source Code},label={code:simcheck}]
import numpy as np
from datetime import datetime, timezone
from typing import Any

def sim_check(sim_data: Any) -> None:
    """
    Validates and prints a diagnostic summary of the simulation data structure.

    Checks key properties of the simulation state (such as start time, step time,
    pointing spheres, celestial bodies, satellites, and detector configurations)
    and reports warnings if crucial simulation elements are missing or uninitialized.

    Args:
        sim_data: The main simulation data structure (SimulationState).
    """
    print("========================================")
    print("--- Simulation Data Check ---")

    # Check for start_time
    if hasattr(sim_data, 'start_time') and sim_data.start_time is not None:
        print(f"          Start Time: {sim_data.start_time}")
    else:
        print("*** Start Time: Not found")

    # Check for delta_time
    if hasattr(sim_data, 'delta_time') and sim_data.delta_time is not None:
        print(f"          Delta Time: {sim_data.delta_time}")
    else:
        print("*** Delta Time: Not found")

    # Check for pointing_spheres
    if hasattr(sim_data, 'pointing_spheres') and sim_data.pointing_spheres:
        print(f"          Pointing Sphere Keys: {list(sim_data.pointing_spheres.keys())}")
    else:
        print("*** Pointing Spheres: Not found or empty")

    # Check for celestial
    if hasattr(sim_data, 'celestial') and sim_data.celestial:
        print(f"          Celestial Dictionary Keys: {list(sim_data.celestial.keys())}")
    else:
        print("*** Celestial Dictionary: Not found or empty")

    # Check for satellites
    if (hasattr(sim_data, 'satellites') and sim_data.satellites and
            hasattr(sim_data, 'counts') and sim_data.counts and
            hasattr(sim_data.counts, 'satellites') and sim_data.counts.satellites):
        num_sats = sim_data.counts.satellites
        print(f"          Number of Satellites: {num_sats}")

        if hasattr(sim_data.satellites, 'position') and sim_data.satellites.position is not None:
            if np.all(sim_data.satellites.position == 0):
                print("*** Satellite positions are all zero.")
            else:
                print("          Satellite positions are not all zero.")
        else:
            print("*** Satellite positions not found.")

        if hasattr(sim_data, 'detector') and sim_data.detector:
            detector = sim_data.detector
            for i in range(num_sats):
                # Check a few attributes to see if they are non-zero
                if (hasattr(detector, 'apertureArea') and detector.apertureArea[i] != 0 or
                    hasattr(detector, 'pixelOmega') and detector.pixelOmega[i] != 0 or
                    hasattr(detector, 'qe') and detector.qe[i] != 0):
                    print(f"            - Satellite {i} has a detector.")
                else:
                    print(f"***         - Satellite {i} does not have a detector.")
        else:
            print("*** Detector information not found for satellites.")
    else:
        print("*** Satellites: Not found")

    print("--- End of Check ---")

if __name__ == '__main__':
    # Create a dummy sim_data for demonstration
    from detector import makeBlankDetector
    from minimalsimulation import (
        create_empty_simulation,
        SatellitesState,
        CelestialState,
    )
    dummy_detector = makeBlankDetector(2)
    dummy_detector.apertureArea = np.array([0.785, 0.])
    dummy_detector.pixelOmega = np.array([1e-10, 0.])
    dummy_detector.qe = np.array([0.5, 0.5])

    dummy_sim_data = create_empty_simulation(
        start_time=datetime(2025, 1, 1, 0, 0, 0, tzinfo=timezone.utc),
        delta_time=60.0
    )
    dummy_sim_data.counts.satellites = 2
    dummy_sim_data.counts.celestial = 2
    dummy_sim_data.pointing_spheres = {
        'sphere1': np.array([[1, 0, 0], [0, 1, 0], [0, 0, 1]]),
    }
    dummy_sim_data.celestial = CelestialState(
        position=np.zeros((2, 3))
    )
    dummy_sim_data.satellites = SatellitesState(
        position=np.zeros((2, 3))
    )
    dummy_sim_data.detector = dummy_detector
    sim_check(dummy_sim_data)

    print("\n--- Testing with a minimal structure ---")
    minimal_sim_data = create_empty_simulation(
        start_time=datetime(2025, 1, 1, 0, 0, 0, tzinfo=timezone.utc)
    )
    sim_check(minimal_sim_data)


\end{lstlisting}



\chapter{all\_demos.py: Demonstration Suite Runner}
\label{sec:alldemos}
\index{all\_demos.py}

\section*{Overview}
Main entry point to execute themed demo suites ('orbits', 'pointing', 'radiometry').

\section*{Literate Source Code}
The source code of \texttt{all\_demos.py} is documented below. Section label: \ref{sec:alldemos}.

\begin{lstlisting}[language=Python,caption={all\_demos.py Source Code},label={code:alldemos}]
import plotly.graph_objects as go
from typing import Union, List, Callable, Optional

# Import all demo functions
from demo1 import demo1
from demo2 import demo2
from demo3 import demo3
from demogeo import demogeo
from demo_fixedpoints import demo_fixedpoints
from demo_exclusion_table import demo_exclusion_table
from demo_pointing_plot import demo_pointing_plot
from radiometry_test import demoFixed

from demo_lambertian import demo_lambertian
from demo_sky_scan import demo_sky_scan
from demo_pointing_vectors import demo_pointing_vectors
from demo_pointing_sequence import demo_pointing_sequence
from demo_constellation import demo_constellation
from demo_show_geo_search import demo_show_geo_search
from demo_requiredIntegrationTime import demo_requiredIntegrationTime
from fibonacciSearch import test_vector_resorting
from pointing import demo_exclusion_pointing
from demo_gap_time_histogram import demo_gap_time_histogram
from demo_observatories_only import demo_observatories_only


def demo_vector_resorting_plot() -> go.Figure:
    """
    Runs the test_vector_resorting function and returns its figure.
    """
    print("\n--- Starting Demo: Vector Resorting ---")
    return test_vector_resorting()


DEMO_GROUPS = {
    'orbits': [
        ('demo1', demo1),
        ('demo2', demo2),
        ('demo3', demo3),
        ('demogeo', demogeo),
        ('demo_constellation', demo_constellation),
        ('demo_observatories_only', demo_observatories_only),
    ],
    'pointing': [
        ('demo_pointing_plot', demo_pointing_plot),
        ('demo_sky_scan', demo_sky_scan),
        ('demo_pointing_vectors', demo_pointing_vectors),
        ('demo_pointing_sequence', demo_pointing_sequence),
        ('demo_show_geo_search', demo_show_geo_search),
        ('demo_exclusion_pointing', demo_exclusion_pointing),
    ],
    'radiometry': [
        ('demo_fixedpoints', demo_fixedpoints),
        ('demo_exclusion_table', demo_exclusion_table),
        ('demo_lambertian', demo_lambertian),
        ('demo_requiredIntegrationTime', demo_requiredIntegrationTime),
        ('demo_vector_resorting_plot', demo_vector_resorting_plot),
        ('demoFixed', demoFixed),
        ('demo_gap_time_histogram', demo_gap_time_histogram),
    ],
}


def list_demo_groups() -> None:
    """Lists available demo groups and their included demo functions."""
    print("=== VibeVolts Demo Groups ===")
    for group_name, demo_list in DEMO_GROUPS.items():
        print(f"\nGroup: '{group_name}' ({len(demo_list)} demos)")
        for name, func in demo_list:
            print(f"  - {name}")
    print("\nUse run_all_demos(group='<name>') to run a specific group inline.")


def run_all_demos(
    group: str = 'all',
    demos: Optional[List[Union[str, Callable]]] = None,
    save_html: bool = False
) -> None:
    """
    Runs demo functions by suite group, specific list, or all demos.

    Args:
        group (str): Demo group to run: 'orbits', 'pointing', 'radiometry', or 'all'.
        demos (list): Optional list of demo names or functions to run.
        save_html (bool): If True, saves plots to HTML file.
    """
    funcs_to_run = []

    if demos is not None:
        all_name_map = {}
        for g_list in DEMO_GROUPS.values():
            for name, func in g_list:
                all_name_map[name] = func
        for item in demos:
            if callable(item):
                funcs_to_run.append((item.__name__, item))
            elif isinstance(item, str) and item in all_name_map:
                funcs_to_run.append((item, all_name_map[item]))
            else:
                print(f"Warning: Demo '{item}' not recognized. Skipping.")
    elif group in DEMO_GROUPS:
        funcs_to_run = DEMO_GROUPS[group]
    elif group == 'all':
        for g_list in DEMO_GROUPS.values():
            funcs_to_run.extend(g_list)
        if not save_html:
            print("Notice: Displaying all demos inline can exceed browser WebGL limits (16 max).")
            print("Defaulting to save_html=True. Set group='orbits', 'pointing', or 'radiometry' for inline notebook display.")
            save_html = True
    else:
        print(f"Unknown group '{group}'. Available groups: {list(DEMO_GROUPS.keys())} or 'all'.")
        return

    figs = []
    print(f"--- Running Demos (Group: {group}, Count: {len(funcs_to_run)}) ---")
    for name, func in funcs_to_run:
        print(f"\n\033[91m--- Executing {name} ---\033[0m")
        res = func()
        if isinstance(res, go.Figure):
            figs.append(res)
        elif isinstance(res, tuple):
            for item in res:
                if isinstance(item, go.Figure):
                    figs.append(item)

    if save_html:
        filename = "all_demo_plots.html" if group == 'all' else f"demo_plots_{group}.html"
        with open(filename, "w") as f:
            f.write("<html><head><title>VibeVolts Demos</title></head><body>")
            f.write(f"<h1>VibeVolts Demo Plots ({group.capitalize()})</h1>")
            for i, fig in enumerate(figs):
                title = fig.layout.title.text if fig.layout.title else f"Plot {i+1}"
                f.write(f"<h2>{title}</h2>")
                f.write(fig.to_html(full_html=False, include_plotlyjs='cdn'))
            f.write("</body></html>")
        print(f"\n--- Demos complete. {len(figs)} plots saved to {filename} ---")
    else:
        print(f"\n--- Displaying {len(figs)} demo plots inline ---")
        for fig in figs:
            fig.show()


if __name__ == '__main__':
    run_all_demos()

\end{lstlisting}



\printindex

\end{document}
